\documentclass[reprint,aps,prx,showpacs,bibfootnote,preprintnumbers,superscriptaddress,amsmath,amssymb]{revtex4-2}
% The EigenEngine house layer on the REVTeX4-2 base. Source of truth:
% templates/eigen-revtex/eigen-revtex.sty (co-located here so the pipeline
% target resolves the .sty at compile time, matching the repo's co-located-sty
% idiom, e.g. eigenengine.sty). It pulls the canonical maths/notation/float
% stack, the shared palette, phys-coloured links and the \claim tag - the
% documentclass carries the REVTeX options. See tex_template.LAYER_PROVIDES.
\usepackage{eigen-revtex}
% The artefact is a NAMED, VERSIONED object (the author's standing law): the
% PDF carries its identity in its own metadata, never shipped as a bare main.pdf.
% pdfsubject/pdfkeywords are revision-NEUTRAL here - the version stamp is applied
% at finalize time (owned by C1), not carried as a stale literal in the source.
\hypersetup{
  pdftitle={Spectral Surface Generation},
  pdfauthor={Gerard McCaul},
  pdfsubject={Reconstructing Gaussian-filtered spectral surfaces from a single overlap-acquisition campaign, with a controlled subset-sum benchmark},
  pdfkeywords={spectral surface generation, Gaussian filtering, quantum spectral estimation, overlap acquisition, subset-sum benchmark, ITQDE},
}

% CONFIDENTIAL marking (internal invention-disclosure white paper).
% A document-STATUS marking, not a scientific disclaimer. ONE macro holds the text.
% To drop it for a future external or journal version, delete this block together
% with the one \notfordistribution banner line after \maketitle. House palette, acc
% then faint, unobtrusive but present on every page. It uses \textemdash and
% \textperiodcentered (commands, so the house punctuation surface stays clean).
\newcommand{\notfordistribution}{%
  \textcolor{acc}{\textbf{CONFIDENTIAL}}\ \textcolor{faint}{\textemdash\ INTERNAL TECHNICAL DOCUMENT \textperiodcentered\ Prepared in support of an invention-disclosure filing \textperiodcentered\ Not for distribution}}
% every-page footer, the marking centred (discreet, faint) with the page number right.
\fancyfoot[C]{\footnotesize\notfordistribution}
\fancyfoot[R]{\small\textcolor{faint}{\thepage}}

\begin{document}

\title{Spectral Surface Generation}
\author{Gerard McCaul\orcidlink{0000-0001-7972-456X}}
% No institutional affiliation on this paper (the author's standing instruction,
% 2026-07-10 - the prior affiliation and its institutional address are scrubbed).
% \noaffiliation is REVTeX's sanctioned "no address" declaration - no placeholder
% institution is invented.
\noaffiliation

\begin{abstract}
Spectral Surface Generation (SSG) separates spectral acquisition from spectral interrogation: one overlap-acquisition campaign, acquired to a stated reach, is reused across spectral location and every supported coarser scale without further quantum evolution. A quantum device samples the real-time overlap signal $G(t)=\langle\psi|e^{-iHt}|\psi\rangle$ once, at fixed Gauss-Hermite grid times, and banks the estimates - a classical layer then reconstructs the Gaussian-filtered spectral surface $Z_\tau(s)=\sum_j p_j\,e^{-\tau(E_j-s)^2}$ and its response channels, the filtered mean, variance and higher cumulants, at a spectral location $s$ and inverse-scale resolution $\tau$ chosen after acquisition. The reuse is exact and bounded: the reconstruction identity is an exact continuous transform, and one banked record supports a finite $(s,\tau)$ envelope set by the record's time reach, node retention and the shot-noise floor. A controlled subset-sum family supplies the validation: at the production weight scale every reconstructed quantity carries an independently computable classical reference, and the construction is checked against it by exact diagonalisation to $n=12$ and matrix-product simulation to $n=20$ under synthetic finite-shot sampling. The benchmark tests reconstruction against that reference rather than asserting a complexity separation.
\end{abstract}

\maketitle

\begin{center}\small\notfordistribution\end{center}

% PLAN:N1
% NODES:
\section{Introduction}\label{sec:N1}

% N1-p1
Questions about a many-body Hamiltonian $H$ probed through a state $|\psi\rangle$ reduce to questions about the probe-weighted spectral measure $\mu_\psi = \sum_j p_j\,\delta(E - E_j)$ - where its weight sits, at which energy it is read, and at what resolution. A quantum device that evolves $|\psi\rangle$ under $H$ accesses this measure through the real-time overlap $G(t) = \langle\psi|e^{-iHt}|\psi\rangle$, its characteristic function, whose frequencies are the spectrum and whose amplitudes are the spectral weights. The operational problem is to turn a finite set of measured overlap values into spectral statements at chosen locations and resolutions, at stated accuracy - to read regularised spectral quantities at a chosen location and scale, not to recover individual eigenvalues.

% N1-p2
Conventional quantum spectral protocols implement the spectral question on the device. Phase-estimation-type readouts, coherently synthesised filters and subspace or eigenstate-targeting schemes fix the target energy and the effective resolution in the circuit that runs, so a new location or a new resolution means a new quantum experiment, and the measurement record produced for one question does not in general serve the next. This is no defect of any particular scheme, and Sec.~\ref{sec:N2} sets the present method against published ones task by task. What the designs share is a single assumption, uniform enough to pass unstated - that spectral location and resolution are parameters of acquisition.

% N1-p3
Spectral Surface Generation (SSG) relaxes this assumption: the spectral filter moves off the device and into classical post-processing. Acquisition still commits to a maximum evolution reach, and that reach fixes the finest supported resolution. What moves after acquisition is selection within that ceiling: a single campaign acquired to a chosen maximum reach supports subsequent selection of spectral location and every coarser resolution within the resulting finite-record envelope. A quantum device samples the spectral signal $G(t)$ at fixed Gauss-Hermite grid times $t_k$ in one acquisition campaign, and the estimates $\hat{G}(t_k)$ are banked. A classical layer then reconstructs the Gaussian-filtered spectral surface at a spectral location $s$ and inverse-scale resolution $\tau$ chosen after acquisition. New $s$, new $\tau$ or a coefficient insertion means only new reconstruction coefficients applied to the stored record - no further quantum evolution. We carry the derivation in the spectral coordinate $s$, the actual spectral location at which the Gaussian window is centred - its relation to the parent ITQDE target convention is developed where that machinery is reconstructed, in Sec.~\ref{sec:N2}. A single property organises the method: the measured overlaps are independent of the target and support a range of resolutions. The acquired record is therefore no longer attached to one spectral question - it becomes a classical interface through which an entire supported family of spectral questions can be asked, each answered by reconstruction alone. Sec.~\ref{sec:N2} proves the independence for the measured objects, and Sec.~\ref{sec:N3} derives the reconstruction identity behind the dials. What one record makes readable, at what accuracy, is bounded in Sec.~\ref{sec:N4}.

% N1-p4
The principal object is the spectral surface $Z_\tau(s) = \sum_j p_j\, e^{-\tau (E_j - s)^2}$ - a Gaussian window of width $1/\sqrt{\tau}$, centred at $s$, applied to the probe-weighted spectral measure - so small $\tau$ is a broad survey and large $\tau$ a narrow local readout. From the same banked record follow the generating surface $\Phi_\tau(s) = \log Z_\tau(s)$, whose derivatives are the response channels - the filtered mean, the filtered variance and the higher cumulants - and, beyond them, within-window ratios and observable branches. Figure~\ref{fig:nf1} fixes the architecture before any formalism: the record is acquired once, and the spectral location $s$ and the resolution $\tau$ enter only as post-acquisition classical dials in reconstruction, leaving the acquired record independent of both. Panel (a) is the pipeline $\hat{G}(t_k) \to c_k(s,\tau) \to \hat{Z}_\tau(s)$, with the quantum-classical boundary drawn once and the three causes of the finite-record envelope labelled at their sources. Panel (b) is a three-level toy measure under three Gaussian kernels, its surface values computed as exact spectral sums, with the envelope causes the panel labels derived in Sec.~\ref{sec:N4}.

% N1-p5
The reuse is bounded, and the bounds are part of the method's definition. Hats mark finite-record objects: $\hat{Z}_{\tau,m}(s)$ is reconstructed from $m$ retained nodes, while the unhatted $Z_\tau(s)$ is reserved for the exact continuous object and the algebraic identities. One banked record supports a finite $(s, \tau)$ domain, the finite-record envelope, bounded by three causes derived in Sec.~\ref{sec:N4}: time reach from the Gaussian aperture, node retention from the Gauss-Hermite gate with its measured retention law $\bar{m} \sim \|H\|\sqrt{\tau}$, and the shot-noise floor at shot budget $K$ per retained estimator. The shot fluctuations propagate through the linear reconstruction with a derived covariance for estimators that share one record. The cost of the classical layer is accounted separately, under an explicitly named query model in Sec.~\ref{sec:N3}. An absolute read of exponentially small surface values is expensive in shots, so the method reads no exponentially small spectral weight at polynomial shot cost, and every later small-weight statement is sized against that covariance and floor.

% N1-p6
We validate the method on a controlled subset-sum family chosen for three diagnostic properties: an exact classical comparator, controlled degeneracy, and a structured route from a commuting limit into a noncommuting regime. The instance $a = (a_1, \ldots, a_n)$ is encoded in the diagonal Hamiltonian $H_0$ and never in the probe. At $V = 0$ the overlap record collapses exactly to computable Fourier coefficients, so every reconstructed quantity has an independently computable classical reference, and at the production weight scale $a_i = O(n^2)$ dynamic programming produces that reference in polynomial time. Bounded noncommuting perturbations $V$ supply the diagnostic regime - a multiplicity fingerprint under a stated forward model, degeneracy cartography and an integer decoder with an erasure rule - constructed in Sec.~\ref{sec:N5}. A planted equal-weight family probes whether readability and growing classical simulation cost can coexist over the tested range, and its interpretation belongs to Sec.~\ref{sec:N7}.

% N1-p7
One scope statement governs the whole paper. The reconstruction, response and cumulant identities, together with the $V = 0$ coefficient collapse, are exact algebra, and the numbers that test them come from exact diagonalisation at $n \le 12$ and TEBD matrix product state evolution at $n \le 20$ under stated shot budgets. The paper makes no asymptotic-hardness or quantum-advantage claim, and the global complexity boundaries - together with the open hardness question raised by the planted family - are stated once, in full, in Sec.~\ref{sec:N7}.

% N1-p8
The argument runs as a dependency chain. Sec.~\ref{sec:N2} reconstructs the parent ITQDE machinery that produces the banked overlap record, proves the record's independence of the target and locates the method task by task against record-based spectral approaches. Sec.~\ref{sec:N3} defines Spectral Surface Generation and derives the reconstruction, response and cumulant identities that make $s$, $\tau$ and coefficient insertions post-acquisition dials. Sec.~\ref{sec:N4} derives the acquisition chain and bounds the supported $(s, \tau)$ domain of one banked record. Sec.~\ref{sec:N5} instantiates the controlled subset-sum embodiment and defines its exact classical comparator. Sec.~\ref{sec:N6} validates the reconstruction, the fingerprint and the decoder against that comparator. Sec.~\ref{sec:N7} bounds what the results establish and states the open hardness question. Sec.~\ref{sec:N8} closes with the three statements the body supports. The appendices carry the complete derivations and the data and code manifest.

% figure environment generated by tex_backend from the figure spec
% spec: fig:nf1 (status: banked)
\begin{figure*}[tb]
    \centering
    \includegraphics[width=\textwidth]{figures/NF1.pdf}
    \caption{Panel (a) is a schematic of the Spectral Surface Generation pipeline: a quantum device samples the real-time overlap record $\hat G(t_k)$ on the Gauss-Hermite grid in one fixed overlap acquisition campaign, and every later step is classical reconstruction of the spectral surface $\hat Z_\tau(s)$ through the coefficients $c_k(s, \tau)$. The quantum-classical boundary is drawn once, and the three causes of the finite-record envelope are labelled at their sources - time reach set by the Gaussian aperture, node retention set by the Gauss-Hermite gate, and the usable region set by shot-noise. Panel (b) shows a generic three-level probe-weighted spectral measure, three Gaussian kernels centred at $s_1$, $s_2$ and $s_3$, and the resulting surface values at those locations on a common quantitative axis, computed inline as exact spectral sums. Together the panels fix the architecture: acquisition is separated from interrogation, and the location $s$ and resolution $\tau$ enter only as post-acquisition classical dials.}
    \label{fig:nf1}
\end{figure*}

% PLAN:N2
% NODES: def:spectral-signal def:spectral-weights def:probe-target-shift def:hamiltonian-shift def:filter-kernel res:itqde-spectral-readout def:staircase def:itqde-unitary-kernel def:retained-overlaps irf:overlap-reconstruction def:quadrature-truncation-budget irf:retained-overlap-budget irf:lambda-only-scalar-phase def:classical-post-processing
\section{Relation to ITQDE and record-based spectral methods}\label{sec:N2}

% N2-p1
Every acquisition in this paper is an estimate of the real-time overlap signal
\begin{widetext}
\[
G(t) \;=\; \langle\psi|\,e^{-iHt}\,|\psi\rangle \;=\; \int e^{-iEt}\,\mathrm{d}\mu_\psi(E) \;=\; \sum_j p_j\,e^{-iE_j t},
\]
\end{widetext}
where $\mu_\psi$ is the spectral measure of the probe state $|\psi\rangle$ and the weights $p_j$ are defined below. The signal is the characteristic function of that measure: its frequencies are the eigenvalues of $H$ and its amplitudes are the weights the probe places on them. The device therefore supplies samples of a classical signal - not an eigenvalue, not a density, not a decision bit. The identity is exact for the continuous signal - the finite-shot estimates taken at grid times are the business of later sections.

% N2-p2
The measure itself is atomic,
\[
\mu_\psi \;=\; \sum_j p_j\,\delta_{E_j}, \qquad p_j \;=\; |\langle E_j|\psi\rangle|^2, \qquad \sum_j p_j = 1,
\]
with weights that are real, nonnegative and normalised - the squared overlaps of the probe state on the Hamiltonian eigenbasis. These are the quantities every filtered object of this paper reweights. They depend on the probe state, so every filtered surface constructed below is probe-weighted rather than a property of the spectrum alone.

% N2-p3
Two coordinates locate everything that follows. The target $T$ is the energy value of interest, fixed by the instance, and the spectral coordinate $s$ marks the location at which the spectrum is interrogated. The parent ITQDE construction \cite{mccaul2025snacks} probes at a displacement $\lambda$ from the target, and the two coordinates are related by
\[
s \;=\; T - \lambda .
\]
Every statement of the method is carried in $s$. The displacement $\lambda$ enters only in this bridge relation, as the coordinate in which the parent machinery was originally written, and no later step works in it.

% N2-p4
The coordinate is carried by scalar shifts of the Hamiltonian,
\[
H_T \;=\; H - T I, \qquad H - s I ,
\]
so in the spectral coordinate the probed operator is the single object $H - sI$. Each is a shift of $H$ by a real scalar multiple of the identity, and such a shift changes only the global phase of real-time evolution - the property that carries the record reuse, proved once the measured record is in hand.

% N2-p5
The Gaussian filter kernel is built from the shifted operator \cite{mccaul2025snacks},
\begin{widetext}
\begin{equation}\label{def:filter-kernel}
e^{-\tau (H - sI)^2}, \qquad Z_\tau(s) \;=\; \langle\psi|\,e^{-\tau (H - sI)^2}\,|\psi\rangle \;=\; \sum_j p_j\,e^{-\tau (E_j - s)^2}.
\end{equation}
\end{widetext}
The kernel is diagonal in the Hamiltonian eigenbasis with entries $e^{-\tau(E_j - s)^2}$, so it concentrates spectral weight within a window of order $1/\sqrt{\tau}$ around $s$, and its probe expectation defines the filtered surface $Z_\tau(s)$ - the Gaussian-smoothed rendering of the spectral measure at the resolution set by $\tau$. The kernel is an analysis object, diagonal in the eigenbasis. No circuit implements it, and no imaginary-time evolution is performed on the state.

% N2-p6
The surface supports a family of response channels, one of which is read as the readout \cite{mccaul2025snacks}. The filtered mean
\begin{equation}\label{def:staircase}
\mu_\tau(s) \;=\; s + \frac{1}{2\tau}\,\partial_s \log Z_\tau(s), \qquad H_\tau(s) \;=\; \mu_\tau(s) - T,
\end{equation}
is the staircase channel of the surface, with $H_\tau(s)$ its target-relative form. As $s$ sweeps across isolated spectral clusters, $\mu_\tau(s)$ holds near the locally dominant cluster energy and steps between clusters, and in the high-$\tau$ regime this produces the staircase surface. The staircase is the most legible of the channels read off $Z_\tau(s)$, though only one among several - the filtered variance and the higher filtered cumulants are derived readouts of the same surface, each following as a response identity of the generating surface.

% N2-p7
The experiment that produces these objects reaches the Gaussian filter without imaginary-time evolution: the filter is written as a Gauss-Hermite quadrature over the imaginary-time axis, and each quadrature node is converted into real-time evolution \cite{leamer2024emulation,mccaul2025snacks}. The circuit element is the unitary kernel
\begin{widetext}
\begin{equation}\label{def:itqde-unitary-kernel}
U \;=\; e^{-i\sqrt{\Delta\tau/2}\,H}, \qquad |\psi_{j+1}\rangle = U\,|\psi_j\rangle, \qquad |\psi_{j-1}\rangle = U^{\dagger}\,|\psi_j\rangle, \qquad |\psi_0\rangle = |\psi\rangle,
\end{equation}
\end{widetext}
whose repeated application generates the bidirectional trajectory from the base state. The device therefore runs ordinary unitary evolutions of a fixed elementary step. Here $\Delta\tau$ is a quadrature step of the classical construction rather than a physical imaginary-time interval - the evolution is unitary real-time throughout, nothing non-unitary is applied to the state, and the imaginary-time character of the filter survives only in the classical weights later attached to the trajectory.

% N2-p8
Along this trajectory the device measures the overlap family that constitutes the banked record. For an observable decomposed as $O = \sum_\ell c_\ell O_\ell$, the retained lists and the reconstruction they support are
\begin{widetext}
\[
\begin{gathered}
A_j \;=\; \langle \psi_{2j} | \psi_{-2j} \rangle, \qquad B_{j,\ell} \;=\; \langle \psi_{2j} |\, O_\ell\, | \psi_{-2j} \rangle,\\[2pt]
\langle N_\tau(s) \rangle \;=\; \sum_{j} \alpha_j(s)\, B_j, \qquad \langle Z_\tau(s) \rangle \;=\; \sum_{j} \beta_j(s)\, A_j, \qquad \langle H_\tau(s) \rangle \;\approx\; \frac{\langle N_\tau(s) \rangle}{\langle Z_\tau(s) \rangle},
\end{gathered}
\]
\end{widetext}
following the parent construction \cite{mccaul2025snacks}, where $B_j = \sum_\ell c_\ell B_{j,\ell}$ gathers the decomposed insertions with their coefficients and the sums run over the retained quadrature nodes. The filtered objects therefore reconstruct from the banked lists as discrete sums, with the filtered ratio their quotient. The quadrature coefficients $\alpha_j$ and $\beta_j$ carry the entire coordinate dependence, so the measured entries never reference the location at which the surface is later evaluated. The decomposition also fixes which numerator overlaps exist in the record - an insertion outside the decomposed family requires new measurements.

% N2-p9
The record is finite because the quadrature is truncated. Keeping $\bar m$ retained overlap pairs discards a tail whose contribution is an additive error
\[
\varepsilon \;\le\; \sum_{k > \bar m} w_k f_k, \qquad \bar m \;\sim\; \|H\|\,\sqrt{\tau \ln(1/\varepsilon_w)},
\]
with $w_k$ the Gaussian quadrature weights, $f_k$ the corresponding integrand values and $\varepsilon_w$ the tolerance on the discarded weight tail. The error decays super-exponentially in $\bar m$ at fixed $\tau\|H\|^2$, because the Gaussian quadrature weights collapse once the grid covers the measured range \cite{leamer2024emulation,mccaul2025snacks}. The $\bar m$-sized retained subset is the device-level measurement load of one acquisition campaign. The square-root budget law and its grid constant follow from the acquisition analysis, superseding the compact budget form given in \cite{mccaul2025snacks}.

% N2-p10
A change of spectral location reaches this record through the scalar shift alone. Under the shift to $H - sI$ the real-time propagator factorises as $e^{-i t (H - sI)} = e^{i t s}\, e^{-i t H}$, so the shift enters only as a global phase \cite{mccaul2025snacks}. The overlap family is measured under the unshifted kernel $U = e^{-i\sqrt{\Delta\tau/2}\,H}$ of Eq.~\eqref{def:itqde-unitary-kernel}, so the phase attaches to the classical coefficients $\alpha_j$ and $\beta_j$ rather than to the measured entries, and the banked family of $A_j$ and $B_{j,\ell}$ is independent of $s$, and of the target $T$. This proves the one narrow claim of the introduction at the level of the measured objects: one record, acquired once, supports the readout at every spectral location. The independence holds for the measured family itself - the accuracy of any reconstruction evaluated from the finite family is a separate finite-record question.

% N2-p11
The independence just proved opens the protocol's cost-free zone. Classical post-processing is any computation on the already-banked overlap family that requires no further quantum experiments,
\begin{widetext}
\[
\{A_j,\, B_{j,\ell}\}\ \text{fixed at acquisition} \;\longrightarrow\; \bigl\{Z_\tau(s),\ \mu_\tau(s),\ \sigma_\tau^2(s)\bigr\}\ \text{at every } (s, \tau),
\]
\end{widetext}
and it includes the sweep of the spectral location $s$, the sweep of the resolution $\tau$, the assembly of response channels such as the staircase, and the later benchmark analyses. The reconstruction applies coordinate-dependent coefficients to the stored record, so the quantum cost is fixed at acquisition and every later evaluation is a classical reweighting. Cost-free refers to quantum experiments alone - every classical evaluation remains bounded by the noise floor of the finite banked record.

% N2-p12
Assembled, this section delivers one object to the reconstruction: the target-independent retained overlap family, evaluated from the same record at every spectral location. The Gaussian-filtered surface $Z_\tau(s)$ of Eq.~\eqref{def:filter-kernel}, the filtered mean of Eq.~\eqref{def:staircase} and the filtered variance are the readouts that Sec.~\ref{sec:N3} reconstructs from that family and owns as its contract, stated there at the level of the continuous quadrature, with truncation and shot noise entering through the finite-record analysis of the acquisition section. What this section establishes is the primitive and its device-level independence of $s$, so the sweep over spectral locations is absorbed entirely into classical post-processing.

% N2-p13
The assembled readout can now be located against its neighbours, beginning with the methods that invert the same time signal for its poles. Pole recovery treats samples of $G(t)$ as a finite sum of complex exponentials and inverts for the frequencies, from Prony's construction \cite{prony1795essai} through the matrix pencil \cite{hua1990pencil} and ESPRIT \cite{roy1989esprit} estimators, and quantum eigenvalue estimation by phase estimation \cite{obrien2019qpe} or by time-series analysis of $G(t)$ \cite{somma2019eigenvalue} targets the same poles. The inversion step is where poor conditioning concentrates when poles are close or data are noisy. The present method performs no inversion for poles at any point. It evaluates regularised linear functionals of the spectral measure at a chosen scale - a linear map applied to the record - so no eigenvalue estimate is produced or claimed, and the conditioning theory of exponential inversion is neither inherited nor needed.

% N2-p14
A closer family shares not just the signal but the measured primitive. Real-time subspace methods - variational quantum phase estimation \cite{klymko2022vqpe}, quantum filter diagonalisation \cite{parrish2019qfd}, real-time Krylov diagonalisation \cite{shen2023realtimekrylov} and fast-forwarding constructions \cite{cortes2022fastforward} - measure overlaps of exactly the kind banked here, but assemble them into overlap and Hamiltonian matrices on a time-evolved basis and solve a generalised eigenvalue problem for eigenpairs. The retained family in this method is never assembled into an eigenproblem. It is reweighted by quadrature coefficients. The acquisition primitive is shared while the task and the post-processing are not, and the comparison concerns exactly those.

% N2-p15
A third family is organised around the filter itself. On the classical side, the kernel polynomial method \cite{weisse2006kpm} banks Chebyshev moments of the spectral density once and applies damping kernels classically - the closest architectural relative - and Heisenberg-limited estimation \cite{lin2022heisenberg} reconstructs the integrated spectral measure from time samples by classical Fourier weighting. On the quantum side, coherent filter synthesis - ground-state filtering \cite{ge2019filter}, quantum eigenvalue transformation of unitaries \cite{dong2022qetu} and the quantum singular value transformation \cite{martyn2021grandunification} - realises the filter as a circuit, so each target energy and each resolution incurs the quantum resource cost again. Here the Gaussian filter of Eq.~\eqref{def:filter-kernel} is applied entirely in post-processing. Moving the filter from the device to the record is the architectural point, since the moment identities for smoothed measures are not themselves new - the claim concerns where the filter is applied.

% N2-p16
The record itself has a class. Classical shadows \cite{huang2020predicting} bank randomised measurement outcomes once and reuse them for many observables, the randomised-measurement toolbox \cite{elben2020crossplatform} reuses one such record across devices, and shadow spectroscopy \cite{chan2025shadow} banks the same randomised records for spectral estimation, while randomised benchmarking \cite{knill2008rb} and cross-entropy benchmarking \cite{gao2024xeblimits} bank sequence and sample statistics reused for many figures of merit, and peaked-circuit constructions \cite{aaronson2024peaked} likewise separate a banked acquisition from later verification analyses. The retained overlap family belongs to this measure-once class, with a difference of object: it is a coherent time record of one fixed probe evolving under the instance Hamiltonian, and its reuse ranges over the spectral location, the resolution and the decomposed observable insertions rather than over arbitrary states or observables. The reuse is probe-fixed - a new probe state, or an observable outside the decomposed family, requires a new acquisition campaign. The contribution of the method is this acquisition-and-reuse architecture together with its finite-record analysis.

% PLAN:N3
% NODES: def:protocol-output-specification def:classical-dial def:spectral-surfacing irf:tau-is-a-classical-dial def:spectral-coordinate-s def:probe-target-shift def:spectral-surface def:spectral-weights def:filter-kernel def:generating-surface thm:fourier-reconstruction-identity def:spectral-signal def:coefficient-update def:retained-overlaps def:filtered-partition def:filtered-mean def:filtered-variance def:centered-readout thm:staircase-cumulant-identities def:staircase thm:response-identities-s res:tau-direction-response thm:natural-parameters-cumulants irf:moment-identities ds:insertion-inflates-coefficient-norm ds:variance-amplification-on-kappa-r irf:heat-kernel-scale-space def:surfacing-algorithm ds:uniform-grid-fft-cost ds:scale-family-cost def:observable-surfaces irf:non-hermitian-emulation-boundary irf:non-hermitian-spectral-emulation
\section{Spectral Surface Generation}\label{sec:N3}

% N3-p1
Spectral Surface Generation is fixed by its contract. Its inputs are the Hamiltonian $H$, the probe state, a spectral bandwidth bound, and the record $\widehat G(t_k)$ measured once on the retained grid - the target-independent overlap family banked in the acquisition section (Sec.~\ref{sec:N4}) with truncation budget $\bar m$. Its output is
\begin{equation}\label{def:protocol-output-specification}
\bigl\{\widehat G(t_k)\bigr\}\ \longmapsto\ \bigl\{\widehat Z_{\tau,m}(s),\ \widehat\mu_\tau(s),\ \widehat\sigma_\tau^{2}(s),\ \widehat\kappa_m(\tau,s),\ \dots\bigr\},
\end{equation}
the finite reconstruction of the Gaussian spectral surface over a supported $(s,\tau)$ domain, with its derived response channels - filtered mean, filtered variance, higher cumulants, within-window ratios and optional observable-branch diagnostics - across a polynomial family of post-measurement settings in $s$ and $\tau$. The hat marks finite-record estimators throughout, and the supported domain, bounded by record reach, node retention and the shot-noise floor, is assembled in the acquisition section. The output is a reusable spectral-surface object rather than a single target readout.

% N3-p2
Every entry in that family of settings is reached by one kind of control. A classical dial is a protocol control adjusted after the quantum measurement by changing the reconstruction coefficients applied to the stored overlap bank,
\begin{equation}\label{def:classical-dial}
(s,\ \tau,\ \text{insertion})\ \longmapsto\ \{c_k\},\qquad \{\widehat G_k\}\ \text{fixed},
\end{equation}
with no new quantum experiment. The essential surfacing dials are the spectral location $s$, the inverse scale $\tau$ up to the range the measured record supports, and derived coefficient insertions for moments, cumulants, ratios and diagnostic sweeps. The insertions divide by what the stored record already holds: location derivatives and cumulants are read from the denominator record alone as powers of the node time, commuting observable branches require the corresponding branch signal acquired in advance, and non-commuting diagnostic branches need their own separate records (Sec.~\ref{sec:N4}). The term is deliberate - it separates quantum-sourced resources from classical post-measurement choices, and it claims no freedom beyond the acquired record.

% N3-p3
Spectral Surface Generation turns the parent protocol \cite{mccaul2025snacks} into an interface between Hilbert space and a classical signal,
\begin{widetext}
\begin{equation}\label{def:spectral-surfacing}
G(t)\ \xrightarrow{\ \text{measured once}\ }\ \bigl\{\widehat G(t_k)\bigr\}\ \xrightarrow{\ \text{classical reconstruction}\ }\ Z_\tau(s)\quad \text{over } (s,\tau).
\end{equation}
\end{widetext}
The device samples the spectral characteristic function $G(t)$ once, and the classical layer reconstructs from the banked record a two-parameter Gaussian spectral surface over spectral location and inverse scale, so every filtered spectral readout, with its derivatives and cumulants, follows from one measurement bank by classical post-processing alone. One bank means one acquisition campaign out to a stated time reach, not a single shot, and the shot accounting is the business of the acquisition analysis.

% N3-p4
That claim rests on a grid fact, asserted here and derived in the acquisition analysis. The retained-node grid is uniform in time with a $\tau$-independent spacing $\Delta t=\pi/(\sqrt{8}\,\|H\|)$, a Nyquist-type spacing at the spectral bandwidth - this uniform grid is the Plancherel-Rotach bulk form of the Gauss-Hermite quadrature of the reconstruction identity derived below, whose retained-node count is the source protocol's quadrature-truncation budget \cite{mccaul2025snacks,leamer2024emulation} - and $\tau$ enters the measurement only through the record reach $t_{\max}$~\eqref{res:aperture-resolution}, the largest retained evolution time. One family measured out to $t_{\max}$ therefore serves every $\tau\le t_{\max}^{2}/(4\log(1/\epsilon_w))$, with $\epsilon_w$ the envelope truncation tolerance, by classical reweighting of the same record - per-$\tau$ retention selects a subset of the measured family, and the reach statement uses the aperture convention $e^{-t^{2}/(4\tau)}$ fixed by the reconstruction identity. Both filter parameters are classical dials in the sense of \eqref{def:classical-dial} - the device measures one family, and every point of the surface is a classical read of it via the surfacing map \eqref{def:spectral-surfacing}.

% N3-p5
Both dials are phrased on location and scale, and the location coordinate comes first. The spectral coordinate is
\begin{equation}\label{def:spectral-coordinate-s}
s = T - \lambda,
\end{equation}
the actual spectral position at which the Gaussian window is centred, while the parent probe shift $\lambda$ is a displacement from the chosen target $T$. We carry the derivation in $s$, and $\lambda$ appears only where a statement is translated back to the parent convention - the parent-convention identities are stated on $\lambda$, and the map converts them without loss.

% N3-p6
At the centre of the contract sits one object. The spectral surface is
\begin{equation}\label{def:spectral-surface}
Z_\tau(s)=\int e^{-\tau(E-s)^{2}}\,d\mu_\psi(E)=\sum_j p_j\,e^{-\tau(E_j-s)^{2}},
\end{equation}
the filter kernel of the parent convention written at window centre $s$ and integrated against the spectral measure of the probe state, the sum form following by evaluating the measure on the eigenbasis with weights $p_j=|\langle E_j|\psi\rangle|^{2}$. The window is centred at spectral location $s$ with width $1/\sqrt{\tau}$, so the energy resolution at scale $\tau$ is $\Delta E_\tau\sim 1/\sqrt{\tau}$ - small $\tau$ gives a broad survey, large $\tau$ a narrow local readout. $Z_\tau(s)$ is a normalising generating function for the filtered ensemble and not a Gibbs partition function, so no thermodynamic reading is implied.

% N3-p7
The object every response reads from is the logarithm of the spectral surface \eqref{def:spectral-surface}. The generating surface is
\begin{equation}\label{def:generating-surface}
\Phi_\tau(s)=\log Z_\tau(s),
\end{equation}
whose derivatives in $s$ and $\tau$ are the response channels of the filtered ensemble, and Spectral Surface Generation is the map from the sampled signal $G(t)$ to $\Phi_\tau(s)$ - from a Hilbert-space signal to a spectral scale surface over location and inverse scale. The continuous surface is strictly positive for $\tau>0$, so the generating surface is defined at every point of the plane, and finite-record caveats enter only with the discrete reconstruction.

% N3-p8
The reconstruction identity connects both objects to the measured signal. The spectral signal $G(t)=\langle\psi|e^{-iHt}|\psi\rangle=\sum_j p_j\,e^{-iE_j t}$ is the characteristic function of the probe-weighted spectral measure \cite{mccaul2025snacks}, and the retained overlap family $\{A_j, B_{j,\ell}\}$ of the acquisition section (Sec.~\ref{sec:N4}) is precisely its measured realisation at the retained nodes. Writing the Gaussian window as a Fourier integral, $e^{-\tau(E-s)^{2}}=\tfrac{1}{2\sqrt{\pi\tau}}\int e^{-t^{2}/(4\tau)}\,e^{i(s-E)t}\,dt$, and exchanging the time integral with the spectral sum identifies that sum as $G(t)$, giving
\begin{equation}\label{thm:fourier-reconstruction-identity}
Z_\tau(s)=\frac{1}{2\sqrt{\pi\tau}}\int_{-\infty}^{\infty}e^{-t^{2}/(4\tau)}\,e^{ist}\,G(t)\,dt.
\end{equation}
Both dial claims are visible in the integrand. The spectral location enters only through the phase $e^{ist}$ and the scale only through the Gaussian time aperture $e^{-t^{2}/(4\tau)}$, so measuring $G(t)$ once determines $Z_\tau(s)$ at every $(s,\tau)$ whose aperture the record supports. \emph{Exact} refers to this continuous transform - the finite-record version carries the hat.

% N3-p9
On the grid $t_k=k\Delta t$ with banked estimates $\widehat G_k$ and node budget $m$ fixed by the truncation budget of the acquisition section (Sec.~\ref{sec:N4}), discretising the identity \eqref{thm:fourier-reconstruction-identity} by Gauss-Hermite quadrature of its Gaussian aperture - whose retained nodes are, in the Plancherel-Rotach bulk, uniformly spaced by $\Delta t$ - gives
\begin{widetext}
\begin{equation}\label{def:coefficient-update}
\widehat Z_{\tau,m}(s)=\sum_{k=-m}^{m}c_k(\tau,s)\,\widehat G_k,\qquad c_k(\tau,s)=\frac{\Delta t}{2\sqrt{\pi\tau}}\,e^{-t_k^{2}/(4\tau)}\,e^{ist_k}.
\end{equation}
\end{widetext}
This is where the contract's dials become operational. The coefficient magnitude $\Delta t\,e^{-t_k^{2}/(4\tau)}/(2\sqrt{\pi\tau})$ is the bulk Gauss-Hermite quadrature weight of the Gaussian aperture, so \eqref{def:coefficient-update} is that quadrature written in Plancherel-Rotach form rather than a uniformly-weighted Riemann sum, and the favourable retained-node count is inherited from the source protocol \cite{mccaul2025snacks,leamer2024emulation}. A retained bank is a sampled characteristic function rather than a list of single-purpose overlaps \cite{somma2019eigenvalue,obrien2019qpe}, and changing $s$, changing $\tau$ or inserting a derived factor is a coefficient update applied to the stored record, with no further quantum evolution. $\widehat Z_{\tau,m}$ is notationally distinct from the continuous $Z_\tau$ - \emph{exact} is reserved for the continuous transform and the algebraic identities and never describes the finite sum - and the two-sided sum is realised from a one-sided record by conjugate completion, established in the acquisition analysis.

% N3-p10
The identities inherited from the parent protocol are stated on the probe dial. In that convention the filter kernel defines the filtered partition function,
\begin{equation}\label{def:filtered-partition}
Z_\tau(\lambda)=\langle\psi|e^{-\tau(H-T+\lambda)^{2}}|\psi\rangle=\sum_j p_j\,e^{-\tau(E_j-T+\lambda)^{2}},
\end{equation}
the normalising partition of the spectral weights, whose kernel-weighted first moment is the filtered mean,
\begin{widetext}
\begin{equation}\label{def:filtered-mean}
\mu_\tau(\lambda)=\frac{\langle\psi|H\,e^{-\tau(H-T+\lambda)^{2}}|\psi\rangle}{Z_\tau(\lambda)}=\frac{1}{Z_\tau(\lambda)}\sum_j p_j\,E_j\,e^{-\tau(E_j-T+\lambda)^{2}},
\end{equation}
\end{widetext}
the kernel-weighted average of the eigenvalues, and whose second moment about that average is the filtered variance,
\begin{equation}\label{def:filtered-variance}
\sigma_\tau^{2}(\lambda)=\frac{1}{Z_\tau(\lambda)}\sum_j p_j\,\bigl(E_j-\mu_\tau(\lambda)\bigr)^{2}\,e^{-\tau(E_j-T+\lambda)^{2}},
\end{equation}
completed by the centred readout,
\begin{equation}\label{def:centered-readout}
m_\tau(\lambda)=\mu_\tau(\lambda)-T+\lambda,
\end{equation}
the target-relative readout that tends to zero when the window centres on the target cluster. Under the spectral coordinate $s=T-\lambda$ \eqref{def:spectral-coordinate-s} these are the surface objects at window centre $s$, with $m_\tau=\mu_\tau-s$. $Z_\tau(\lambda)$ and $Z_\tau(s)$ denote one object in two conventions, not two different surfaces.

% N3-p11
For the target-shifted Hamiltonian of the parent convention, write $H_\tau(\lambda)=\mu_\tau(\lambda)-T$ for the staircase, the most legible of the derived response channels. Differentiating $\ln Z_\tau(\lambda)$ of the filtered partition \eqref{def:filtered-partition} under the sum, and differentiating once more for the slope, gives
\begin{widetext}
\begin{equation}\label{thm:staircase-cumulant-identities}
H_\tau(\lambda)+\lambda=-\frac{1}{2\tau}\,\partial_\lambda \ln Z_\tau(\lambda)=m_\tau(\lambda),\qquad \partial_\lambda H_\tau(\lambda)=-2\tau\,\sigma_\tau^{2}(\lambda),
\end{equation}
\end{widetext}
exactly, for every $\lambda$, $\tau$ and coupling. $\ln Z_\tau(\lambda)$ is therefore a cumulant-generating function along the probe dial. The staircase - equal to the centred readout \eqref{def:centered-readout} by the identity above - the filtered variance and every higher filtered cumulant are read from the single measured family by inserting classical factors $(-it_k)^{m}$ into the reconstruction coefficients \eqref{def:coefficient-update}. That factor is the power that differentiating the single-time reconstruction phase $e^{ist_k}$ of \eqref{thm:fourier-reconstruction-identity} in the probe dial brings down. The acquisition and the reference implementation bank the overlap at the doubled node time $2t_k$ of the bidirectional trajectory (Sec.~\ref{sec:N4}), whose phase $e^{-2i\lambda t_k}$ carries the same insertion as $(-2it_k)^{m}$ in that convention. This does not change what the device must record. The $H$-weighted overlap channel is redundant for staircase readout, and the correlation between numerator and denominator shot noise that ratio cancellation requires is native to the log-derivative estimator rather than modelled. The identities are exact algebra for the continuous objects - their finite-record versions carry the hat and inherit the truncation budget of the bank.

% N3-p12
Carried into the spectral coordinate by the chain rule $\partial_\lambda=-\partial_s$, each probe-dial identity \eqref{thm:staircase-cumulant-identities} maps onto its location form, and the two conventions carry the same content \cite{mccaul2025snacks}. Differentiating the generating surface \eqref{def:generating-surface} in $s$ under the sum gives
\begin{widetext}
\begin{equation}\label{thm:response-identities-s}
\partial_s\log Z_\tau(s)=2\tau\bigl(\mu_\tau(s)-s\bigr),\qquad \partial_s^{2}\log Z_\tau(s)=4\tau^{2}\,\sigma_\tau^{2}(s)-2\tau,
\end{equation}
\end{widetext}
so $\mu_\tau(s)=s+\tfrac{1}{2\tau}\,\partial_s\log Z_\tau(s)$ and, differentiating once more, $\partial_s\mu_\tau=2\tau\,\sigma_\tau^{2}(s)$. The moments are taken under the filtered weights $q_j(\tau,s)=p_j\,e^{-\tau(E_j-s)^{2}}/Z_\tau(s)$, so the filtered mean \eqref{def:filtered-mean} and variance \eqref{def:filtered-variance} are read directly off the surface by differentiation alone. They depend on both dials - they are filtered moments under $q_j$, not spectral moments of the bare weights $p_j$.

% N3-p13
The scale direction is a response channel of the same surface. Differentiating in $\tau$ under the sum gives the filtered second moment about $s$, and substituting the first response identity \eqref{thm:response-identities-s} isolates the variance:
\begin{widetext}
\begin{equation}\label{res:tau-direction-response}
-\partial_\tau\log Z_\tau(s)=\sigma_\tau^{2}(s)+\bigl(\mu_\tau(s)-s\bigr)^{2},\qquad \sigma_\tau^{2}(s)=-\partial_\tau\log Z_\tau(s)-\frac{1}{4\tau^{2}}\bigl(\partial_s\log Z_\tau(s)\bigr)^{2}.
\end{equation}
\end{widetext}
The scale is therefore a second response direction of the one generating surface, not only a family parameter to scan - the parent account treated $\tau$ as a setting \cite{mccaul2025snacks}, and the scale response is stated as an addition to the inherited identities. It adds no new measured quantity, being another derivative of the same reconstructed surface within the acquired reach.

% N3-p14
The two response directions - the location identities \eqref{thm:response-identities-s} and the scale response \eqref{res:tau-direction-response} - organise into a single cumulant structure once the square is completed. With the tilt $h=2\tau s$,
\begin{widetext}
\begin{equation}\label{thm:natural-parameters-cumulants}
Z_\tau(s)=e^{-\tau s^{2}}\,\widetilde Z(\tau,h),\qquad \widetilde Z(\tau,h)=\sum_j p_j\,e^{-\tau E_j^{2}+hE_j},\qquad \partial_h^{\,m}\log\widetilde Z\,\Big|_{h=2\tau s}=\kappa_m(\tau,s),
\end{equation}
\end{widetext}
and at fixed $\tau$ the function $\log\widetilde Z$ is a cumulant-generating function in $h$ whose $m$-th derivative at $h=2\tau s$ is the filtered cumulant $\kappa_m(\tau,s)$ of the weights $q_j$. The first cumulant is $\kappa_1=\mu_\tau(s)$ exactly, and not $\mu_\tau(s)-s$, because $\log\widetilde Z=\log Z_\tau+\tau s^{2}$ absorbs precisely the shift that the $s$-derivative of $\log Z_\tau$ carries. Two correct-looking computations disagree here - differentiating $\log Z_\tau$ in the tilt yields $\mu_\tau(s)-s$, both computations are correct, and they name different derivatives, the cumulant reading belonging to $\log\widetilde Z$. The higher cumulants are central and unchanged by the location, $\kappa_2=\sigma_\tau^{2}(s)$ and $\kappa_3$ the filtered third central moment. $(s,\tau)$ are the geometric coordinates of the surface and $(h,\tau)$ the natural ones - moving the window centre is a linear tilt inside a fixed quadratic penalty.

% N3-p15
The cumulant hierarchy \eqref{thm:natural-parameters-cumulants} is algebraically free on the retained bank and carries a statistical cost. Each order inserts one more power of the node time into the coefficients, upweighting the long-time tail where the aperture is thin. Treated as a Gaussian measure in $t$ of width $\sqrt{2\tau}$, the squared coefficients give a norm growth of order $(2\tau)^{m/2}\sqrt{(2m-1)!!}$, so the variance of the reconstructed $m$-th cumulant - a linear estimator in the noisy samples - grows as $(2\tau)^{m}(2m-1)!!$ relative to the per-node noise. Holding fixed accuracy at higher order requires a super-exponentially larger shot budget or stronger regularisation, while the quantum side is unchanged - no new observable and no additional overlap channel. The cost sits in the insertion route - differencing a noisy surface is ill conditioned, amplifying noise as $(\delta s)^{-m}$.

% N3-p16
The scale dial carries the same asymmetry in a sharper analytic form. With $u=1/(4\tau)$, the normalised smoothed spectral density satisfies the heat equation in the distributional sense,
\begin{widetext}
\begin{equation}\label{irf:heat-kernel-scale-space}
\partial_u\rho_u(s)=\partial_s^{2}\rho_u(s),\qquad u=\frac{1}{4\tau},\qquad \rho_u(s)=\sqrt{\tau/\pi}\,Z_\tau(s).
\end{equation}
\end{widetext}
Increasing $u$ - coarsening the scale - is the heat semigroup, a contractive smoothing that damps each spectral frequency $\omega$ by $e^{-u\omega^{2}}$, so coarsening within the acquired record is stable. Decreasing $u$ would multiply each frequency by $e^{+u\omega^{2}}$, the ill-posed backward heat flow, so refinement is never computed by inversion. Refinement is instead a resource statement. Recovering scale $\tau$ requires a record reach $t_{\max}(\tau)=2\sqrt{\tau\log(1/\epsilon_w)}$~\eqref{res:aperture-resolution} and per-node accuracy that grows with $\tau$ through the reconstruction coefficient norm, the smallest recoverable scale is set by the acquired reach with $\tau_{\max}\sim t_{\max}^{2}/(4\log(1/\epsilon_w))$, and the read is floored where the surface drops below the per-node shot-noise level propagated through the coefficient norm. Scale space~\cite{witkin1983scale,koenderink1984structure,lindeberg1994scale} is here a forward-stability statement and never a reversibility one - coarsening is post-processing, while refinement demands longer and more accurate real-time data.

% N3-p17
The analytic asymmetry between location and scale \eqref{irf:heat-kernel-scale-space} reappears as arithmetic once the cost of the classical layer is accounted under an explicit query model, in three output modes. An oracle evaluates $\widehat Z_{\tau,m}(s)$ per $(s,\tau)$ query at cost of order $m$. On a uniform $s$ grid of size $M$ at fixed $\tau$, the apertured bank $g_k=e^{-t_k^{2}/(4\tau)}\,\widehat G_k$ makes the whole $s$ line one discrete Fourier transform, at order $mM$ directly and order $(m+M)\log(m+M)$ by FFT on a commensurate zero-padded grid, the Gaussian aperture being a fixed pre-multiply that adds no order. A scale family at $L$ values of $\tau$ costs $L$ single-scale transforms, because only the phase carries $s$ while the aperture carries $\tau$ and changes the summand itself. These figures cover reconstruction from an already banked record, not acquisition, whose shot costs the acquisition analysis derives.

% N3-p18
The output family \eqref{def:protocol-output-specification} closes with the observable branches, which reuse the coefficient machinery unchanged. For an observable $O$ diagonal in the spectral basis, the numerator surface is
\begin{widetext}
\begin{equation}\label{def:observable-surfaces}
N_{\tau,O}(s)=\sum_j p_j\,O_j\,e^{-\tau(E_j-s)^{2}}=\frac{1}{2\sqrt{\pi\tau}}\int_{-\infty}^{\infty}e^{-t^{2}/(4\tau)}\,e^{ist}\,G_O(t)\,dt,\qquad \langle O\rangle_{\tau,s}=\frac{N_{\tau,O}(s)}{Z_\tau(s)},
\end{equation}
\end{widetext}
where $G_O(t)$ is the observable-inserted overlap signal banked as the $B$ channel of the acquisition section (Sec.~\ref{sec:N4}) \cite{knill1998power}. The derivation repeats the reconstruction identity with $O_j$ inserted, the filtered observable is the ratio, and the denominator is the spectral surface itself. Inserted branches are therefore observable surfaces \eqref{def:observable-surfaces} over the same $(s,\tau)$ domain, reconstructed by the same coefficients applied to the branch record \cite{chan2025shadow}, the earlier reconstruction identity \eqref{thm:fourier-reconstruction-identity} reused rather than a new one derived. The identity as stated covers observables diagonal in the spectral basis - non-commuting diagnostic branches are separate records, costed separately in the acquisition analysis.

% N3-p19
One boundary remains to mark. Reweighting the banked record by other admissible time-domain kernels within the acquired support reconstructs other spectral functions $F(H)$ of the Hermitian generator, and this reweighting freedom is the stated boundary of the Gaussian case - non-Gaussian kernels inherit the same reach and noise floors, and this manuscript treats only the Gaussian case. What the record does not determine is the dynamics of an arbitrary non-Hermitian generator,
\begin{equation}\label{irf:non-hermitian-emulation-boundary}
K=A-iB,\qquad e^{-iKt}=e^{-iAt-Bt},\qquad [A,B]\neq 0,
\end{equation}
whose evolution carries operator-ordering structure absent from the single spectral signal of $H$. Simulating such dynamics \eqref{irf:non-hermitian-emulation-boundary} needs a larger construction - a dilation, a linear-combination-of-unitaries synthesis or additional branch measurements \cite{leamer2024emulation,kamakari2022digital}. The supported claim is spectral non-Hermitian emulation, filtered non-unitary spectral operators from one Hermitian record, and never arbitrary non-Hermitian dynamics.

% PLAN:N4
% NODES: res:ancilla-interferometric-front-end def:iqm-ancilla-readout def:retained-overlaps def:within-window-ratio-observables irf:hermitian-time-reflection def:spectral-signal res:one-sided-reconstruction thm:fourier-reconstruction-identity res:cosine-sine-filter-folding res:half-depth-symmetric-branches def:symmetric-record-projection irf:odd-channel-is-physical def:record-architecture-three-tier res:native-hamiltonian-execution def:native-hardware-V def:shot-budget-K res:retention-frontier-nyquist irf:retained-overlap-budget res:bandwidth-grid def:coefficient-update irf:tau-is-a-classical-dial res:aperture-resolution def:quadrature-truncation-budget res:shot-floor-visibility-wall irf:shared-bank-ratio-boundary
\section{Acquisition and the finite-record envelope}\label{sec:N4}

% N4-p1
The physical layer fixes what one shot can deliver. An ancilla prepared in $|+\rangle$ controls the forward and backward native evolutions of the probe state, $|\psi_{\pm 2j}\rangle = U^{\pm 2j}|\psi\rangle$ with $U = e^{-i\sqrt{\Delta\tau/2}\,H}$, in the Hadamard-test arrangement \cite{knill1998power}. Reading the ancilla $X$ and $Y$ expectations recovers the two quadratures of the denominator overlap,
\[
\langle X\rangle + i\,\langle Y\rangle = \langle\psi_{2j}|\psi_{-2j}\rangle = A_j ,
\]
and inserting a retained observable $O_\ell$ between the two evolutions recovers $B_{j,\ell} = \langle\psi_{2j}|O_\ell|\psi_{-2j}\rangle$ in the same arrangement. Each entry of the retained overlap family is therefore one physical readout, a single-shot dispersive measurement on one dedicated ancilla per shot, with fidelity of at least $0.99$ on current devices \cite{krantz2019engineer, marxer2023long}. The fidelity is a hardware figure quoted from the device literature, and this layer adds no hidden exponential cost above the visibility floor.

% N4-p2
The record must serve a fixed set of targets. For probe pairs $(\lambda_i, \lambda_j)$ lying inside one cluster window of width at most $1/\sqrt{\tau}$, the targets are
\begin{widetext}
\[
R_{ij} = \frac{Z_\tau(\lambda_i)}{Z_\tau(\lambda_j)}, \qquad
\Delta\mu_{ij} = \mu_\tau(\lambda_i) - \mu_\tau(\lambda_j), \qquad
\Delta\sigma^2_{ij} = \sigma^2_\tau(\lambda_i) - \sigma^2_\tau(\lambda_j).
\]
\end{widetext}
Numerator and denominator are estimated from the same retained overlap family, so their shot noise is shared and the variance of the ratio is set by relative rather than absolute error. The observables are defined only for pairs inside a single window, and their error analysis follows from the shared-record covariance.

% N4-p3
These targets require only part of the formally two-sided record. For Hermitian $H$ the spectral measure $\mu_\psi$ is real and positive, so conjugating $G(t) = \int e^{-iEt}\,d\mu_\psi(E)$ only reverses the sign of $t$,
\begin{equation}\label{irf:hermitian-time-reflection}
G(-t) = G(t)^* .
\end{equation}
Negative-time samples carry no independent information, and the minimal independent denominator record is a one-sided complex signal - $G(t)$ on $t \ge 0$, both quadratures on the half-line. The reduction holds for Hermitian generators only, and a measured departure from the conjugation relation is itself a diagnostic channel.

% N4-p4
Applying the conjugation relation~\eqref{irf:hermitian-time-reflection} to the reconstruction identity $Z_\tau(s) = \frac{1}{2\sqrt{\pi\tau}}\int_{-\infty}^{\infty} e^{-t^2/(4\tau)}\,e^{ist}\,G(t)\,dt$ converts the symmetry into an acquisition saving. Splitting the integral at $t = 0$ and conjugating the negative branch combines the two halves into twice a real part, and the identity folds to
\begin{equation}\label{res:one-sided-reconstruction}
Z_\tau(s) = \frac{1}{\sqrt{\pi\tau}}\,\operatorname{Re}\int_0^{\infty} e^{-t^2/(4\tau)}\,e^{ist}\,G(t)\,dt ,
\end{equation}
equivalently $\frac{1}{2\sqrt{\pi\tau}}\big[G(0) + 2\operatorname{Re}\int_0^{\infty} e^{-t^2/(4\tau)}\,e^{ist}\,G(t)\,dt\big]$ once the $t = 0$ endpoint is counted explicitly. The folded form is stated up to that endpoint convention, which must match the discrete coefficient sum the code implements. The negative half of the bank is filled by conjugate completion rather than measured - positive times with both quadratures suffice for every probe $s$, and every later resource count assumes this one-sided record.

% N4-p5
The fold is a property of Hermitian evolution rather than of the Gaussian filter in particular. With $U(t) = e^{-iHt}$, the symmetric and antisymmetric combinations of forward and backward evolution are trigonometric functions of $H$, and any Fourier-representable filter $F(H) = \int \kappa(t)\,U(t)\,dt$ folds onto $t \ge 0$,
\begin{widetext}
\begin{equation}\label{res:cosine-sine-filter-folding}
\begin{aligned}
&\frac{U(t)+U(-t)}{2} = \cos(Ht), \qquad \frac{U(t)-U(-t)}{2i} = -\sin(Ht),\\
&F(H) = \int_0^{\infty}\kappa_S(t)\,\big[U(t)+U(-t)\big]\,dt + \int_0^{\infty}\kappa_A(t)\,\big[U(t)-U(-t)\big]\,dt ,
\end{aligned}
\end{equation}
\end{widetext}
with $\kappa_S(t) = \frac{1}{2}[\kappa(t)+\kappa(-t)]$, $\kappa_A(t) = \frac{1}{2}[\kappa(t)-\kappa(-t)]$, and the sine-channel sign fixed by convention. This places the banked record alongside the cosine-filter constructions of kernel-polynomial and filter-synthesis methods \cite{weisse2006kpm, ge2019filter, dong2022qetu}. The fold concerns filter representation - it does not by itself reduce circuit depth.

% N4-p6
Circuit depth is instead where the symmetric branch implementation acts, not the operator fold~\eqref{res:cosine-sine-filter-folding}. In the branch convention $|\psi(t)\rangle = e^{-iHt}|\psi\rangle$, the symmetric two-branch overlap reaches signal time $T$ with each branch evolving only to depth $T/2$, and the observable branch carries an explicit spectral form,
\begin{widetext}
\begin{equation}\label{res:half-depth-symmetric-branches}
\begin{aligned}
\langle\psi(T/2)|\psi(-T/2)\rangle &= \langle\psi|e^{+iHT}|\psi\rangle = G(T)^* ,\\
G_O(t) = \langle\psi(t/2)|O|\psi(-t/2)\rangle &= \sum_{j,k} c_j^*\,c_k\,O_{jk}\,e^{+i(E_j+E_k)t/2} ,
\end{aligned}
\end{equation}
\end{widetext}
which collapses to $G(-t)$ for $O = I$. The forward-forward object $\langle\psi|e^{-itH/2}\,O\,e^{-itH/2}|\psi\rangle$, with both evolutions of the same sign, is a distinct quantity and must not be equated with the branch observable. The resource consequence sits in an exact class: for a signal at time $T$ the symmetric implementation halves the coherence time and the accumulated phase error per branch, a constant-factor architectural gain at the same asymptotic cost.

% N4-p7
On finite data the symmetry becomes an estimator improvement. We project the measured record onto the Hermitian reflection class and keep the rejected component as a named object,
\begin{widetext}
\begin{equation}\label{def:symmetric-record-projection}
\widehat G_{\mathrm{sym}}(t) = \tfrac{1}{2}\big[\widehat G_+(t) + \widehat G_-(-t)^*\big],
\qquad
\widehat R_{\mathrm{time}}(t) = \tfrac{1}{2}\big[\widehat G_+(t) - \widehat G_-(-t)^*\big].
\end{equation}
\end{widetext}
Under independent unbiased noise of equal variance on the two branches, the variance of $\widehat G_{\mathrm{sym}}$ halves. The residual $\widehat R_{\mathrm{time}}$ should vanish for an ideal Hermitian denominator signal, and a nonzero value is a diagnostic for drift, branch imbalance, time-reversal miscalibration or non-unitary contamination, read as a signal rather than subtracted as a correction. Both objects are hatted finite-record estimators.

% N4-p8
Symmetrisation must not be mistaken for discarding the imaginary quadrature. The odd combination $\frac{1}{2i}[G(t) - G(-t)] = \operatorname{Im} G(t)$ is physical signal, and every probe off the filter centre requires it,
\begin{equation}\label{irf:odd-channel-is-physical}
\operatorname{Re}\big(e^{ist}G(t)\big) = \cos(st)\,\operatorname{Re} G(t) - \sin(st)\,\operatorname{Im} G(t) .
\end{equation}
The projection~\eqref{def:symmetric-record-projection} averages across the pair $+t$ and $-t$, never across the quadratures of $G(t)$, and it removes only the component that violates $G(-t) = G(t)^*$ - the physical antisymmetric part of the quadrature decomposition~\eqref{irf:odd-channel-is-physical} survives intact.

% N4-p9
One acquisition campaign banks three record tiers. The first is the one-sided complex denominator record $G(t)$, with negative times supplied by conjugation and the branch depth $T/2$ of the symmetric implementation~\eqref{res:half-depth-symmetric-branches} - the minimal bank. The second is the commuting moment record for $O = f(H)$, which inherits the same one-sided compression. The third is the set of non-commuting diagnostic branches
\begin{equation}\label{def:record-architecture-three-tier}
B_O(t) = \langle\psi(t)|O|\psi(-t)\rangle = \langle\psi|e^{+iHt}\,O\,e^{+iHt}|\psi\rangle ,
\end{equation}
genuinely two-sided and costed separately. The $A_j$ bank is the minimal surfaced signal, and the $B$ branches are additional channels outside the core spectral surface. The two-sided cost applies to the third tier alone, and the one-sided compression of the first two tiers does not extend to the non-commuting branches.

% N4-p10
The device that generates this record executes the Hamiltonian natively. It evolves under
\[
H = H_0 + V, \qquad V = H_{\mathrm{dev}} - H_0 ,
\]
as its bare Hamiltonian, where native means that $V$ consists of the residual interactions, cross-talk and spectator couplings present without compilation. No Trotter decomposition of $V$ is required \cite{lloyd1996universal}, and the analogue evolution duration $t_j = 2j\sqrt{\Delta\tau/2}$ replaces gate depth as the depth metric. The campaign needs a circuit family of size $O(\bar m M)$ with analogue depth bounded by the longest retained time - a polynomial count under the operationally dense regime with $\bar m$ and $M$ polynomial in $n$, and coherent operation requires that longest time to sit inside the device coherence window. Superconducting, trapped-ion and neutral-atom platforms each supply $V$ natively \cite{yan2018tunable, kamakari2022digital}, and the resulting counts place the $n \le 30$ and $n \approx 50$ regimes against the named classical and quantum subset-sum baselines \cite{hgj2010, bonnetain2020improved}.

% N4-p11
Repetition enters through a single symbol. Let $K$ be the number of independent shots per retained overlap estimator. Each shot returns a binary outcome, so the single-estimator error scales as $K^{-1/2}$ and the additive error of the assembled estimate is
\begin{equation}\label{def:shot-budget-K}
\varepsilon_K = c\,K^{-1/2},
\end{equation}
with $c$ controlled by the retained overlap coefficient norm, which stays polynomial in $n$ when $\bar m$ and $M$ are polynomial in $n$. The bound on $c$ is conditional on exactly those polynomial counts, and the error scale is additive - it bounds absolute rather than relative accuracy.

% N4-p12
Throughout, $\bar m$ is the retained node count of the truncated Gauss-Hermite quadrature introduced with the retained family - the record's device-level measurement load and the source protocol's quadrature-truncation budget \cite{mccaul2025snacks}. Exact diagonalisation measurement across all fifteen $(n, \tau, \mathrm{threshold})$ cells gives $\bar m \approx 3.7\,\|H\|\sqrt{\tau}$, polynomial in $n$ with a fitted log-log slope near $2.5$. The apparent rise at $n = 12$ reflects a single-instance draw in $\|H\|$, shown by the flat collapse column - the frontier is measured on one instance per system size, and instance-to-instance bands are deferred. The same $\|H\|\sqrt{\tau}$ product sets the cost of the classical dynamic-programming lattice \cite{bellman1957dp}, so the law states consistency with the classical baseline rather than advantage. The $\|H\|$ factor comes from the grid and the $\sqrt{\tau}$ factor from the aperture.

% N4-p13
The grid spacing separates into what bandwidth alone dictates and what the step convention ships. The bare argument gives only $\Delta t \lesssim \pi/\|H\|$. The source protocol advances the probe by the elementary step $U = e^{-i\sqrt{\Delta\tau/2}\,H}$ and reads the retained overlaps at the doubled time of the bidirectional trajectory, $\langle\psi|e^{-2it_kH}|\psi\rangle$ at node positions $t_k = k\sqrt{2\,\Delta\tau}$, a signal of spectral bandwidth $2\|H\|$. Matching the physical spacing $\Delta t = \sqrt{2\,\Delta\tau}$ to the Nyquist step at that doubled bandwidth fixes
\begin{equation}\label{res:bandwidth-grid}
\Delta\tau = \frac{\pi^2}{16\,\|H\|^2},
\qquad
\Delta t = \frac{\pi}{\sqrt{8}\,\|H\|} .
\end{equation}
The factor $\sqrt{8} = 2\sqrt{2}$ decomposes exactly: a factor $2$ from sampling the doubled-time overlap $e^{-2itH}$, and a $\sqrt{2}$ oversampling margin inherited from the ITQDE step convention $\sqrt{\Delta\tau/2}$. The spacing derivable from bandwidth alone is $\pi/(2\|H\|)$ - the additional $\sqrt{2}$ is a stated Gauss-Hermite quadrature convention inherited from the source protocol \cite{mccaul2025snacks,leamer2024emulation}, not a Nyquist consequence. The grid is $\tau$-independent, the retained Gauss-Hermite node set in its Plancherel-Rotach bulk, and on it the retained count scales as $t_{\max}/\Delta t$.

% N4-p14
The maximum time follows from the truncation tolerance. Truncating the Gaussian envelope $e^{-t^2/(4\tau)}$ at tolerance $\epsilon_w$ means solving $e^{-t_{\max}^2/(4\tau)} = \epsilon_w$, and the $1/(4\tau)$ in the exponent fixes the factor of two in the root,
\begin{widetext}
\begin{equation}\label{res:aperture-resolution}
t_{\max} = 2\sqrt{\tau\log(1/\epsilon_w)},
\qquad
\bar m = \frac{t_{\max}}{\Delta t} \sim \frac{4\sqrt{2}}{\pi}\,\|H\|\sqrt{\tau\log(1/\epsilon_w)} .
\end{equation}
\end{widetext}
This closes the measured law: the grid supplies the $\|H\|$ factor and the aperture the $\sqrt{\tau}$. The retained count $\bar m$ is the Gauss-Hermite quadrature-truncation budget, small because the Gaussian quadrature is super-exponentially convergent in its order, so the count is set by aperture truncation rather than by resolving the fastest oscillation \cite{mccaul2025snacks,leamer2024emulation}. In resolution language the window width is $\Delta E_\tau \sim 1/\sqrt{\tau}$, so resolving a spectral feature of width $\delta E$ requires $\tau \gtrsim \delta E^{-2}$ and $\bar m \gtrsim \|H\|/\delta E$ - the number of resolvable cells across the band, the aperture-cost correspondence that Heisenberg-limited spectral estimation makes precise \cite{lin2022heisenberg}. The reference numerics compute $t_{\max}$ under the envelope normalisation $e^{-t^2/\tau_{\max}}$, so prose and code differ only in which Gaussian-width convention names $\tau$, and we state the one convention the code obeys after that rescaling.

% N4-p15
Shot noise propagates through the reconstruction linearly, because the reconstruction is a fixed weighted sum of the record - the coefficient-update mechanism $\widehat Z_\tau(s) = \sum_k c_k(\tau,s)\,\widehat G_k$ of Spectral Surface Generation (Sec.~\ref{sec:N3}). With per-node errors $\widehat G_k = G_k + \xi_k$, $\mathbb{E}[\xi_k] = 0$ and $\mathbb{E}|\xi_k|^2 \sim 1/K$ set by the shot budget~\eqref{def:shot-budget-K}, the variance and the shot count for an absolute read at relative accuracy $\epsilon$ follow at once,
\begin{equation}\label{res:shot-floor-visibility-wall}
\operatorname{Var}\widehat Z_\tau(s) \sim \frac{1}{K}\sum_k |c_k|^2 = \frac{C_{\tau,s}^2}{K},
\qquad
K \gtrsim \frac{C_{\tau,s}^2}{\epsilon^2\,Z_\tau(s)^2} ,
\end{equation}
where the coefficient norm is defined by $C_{\tau,s}^2 = \sum_k |c_k(\tau,s)|^2$ and the readout error is $\delta Z_\tau(s) \sim C_{\tau,s}/\sqrt{K}$. This is the visibility wall. Exponentially small $Z_\tau$ is exponentially expensive to read absolutely, a hardness statement about absolute readout rather than a defect of the estimator, and the method reads no exponentially small spectral weight at polynomial shot cost.

% N4-p16
The ratio targets are read from one record, and the shared shots correlate their errors. For the linear reconstruction with stacked real weight vectors $v_i$ and per-node quadrature noise of variance $\sigma^2 \le 1/K$ independent across nodes, two reconstructions $\widehat Z_i$ and $\widehat Z_j$ have covariance $\sigma^2\, v_i \cdot v_j$, and the log-ratio estimator $\widehat u = [\log\widehat Z_i - \log\widehat Z_j]/\Delta$ obeys, in the delta-method regime,
\[
\operatorname{Var}\widehat u
= \frac{\sigma^2}{\Delta^2}\,\bigg\|\frac{v_i}{Z_i} - \frac{v_j}{Z_j}\bigg\|^2
= \frac{\sigma^2}{\Delta^2 P^2}\,\bigg\|\frac{v_i}{\zeta_i} - \frac{v_j}{\zeta_j}\bigg\|^2 ,
\]
where the second form writes $Z_i = P\zeta_i$ with $P$ the window's common absolute weight and $\zeta_i$ an order-one shape factor. The shared record yields a common-mode variance reduction that is a constant in $n$ - large when the two probes are close - while the $P^{-2}$ floor persists. The estimator is defined only where both reconstructions clear the additive shot floor, $Z \ge c\,\sigma\|v\|$ with $c$ of order five, so a window weight of order $g/2^n$ costs exponentially many shots at any fixed accuracy, ratio structure notwithstanding. Every within-window ratio read is conditioned on the window clearing the visibility floor~\eqref{res:shot-floor-visibility-wall}.

% N4-p17
One banked record from one acquisition campaign therefore supports a bounded $(s, \tau)$ region - the finite-record envelope - whose three boundary causes are the maximum evolution time, node retention and the shot-noise floor. Figure~\ref{fig:nf2} draws the envelope. Panel (a) shows the Gaussian aperture over the retained Gauss-Hermite nodes with the Plancherel-Rotach retention threshold, panel (b) shows the supported domain with each boundary cause drawn where it acts, and panel (c) checks the retained-count law against single-instance exact diagonalisation data, with the analytical prediction and the fitted scaling drawn separately. Panels (a) and (b) are schematics, panel (c) is a single-instance check that establishes no instance-to-instance spread, and the figure addresses no reconstruction accuracy below the shot-noise floor. The envelope and the per-shot execution model are the objects on which the cost accounting rests. Table~\ref{tab:resource} collects the independent acquisition resources of one campaign as a plain resource ledger: the time points, the shot budget, the evolution reach, the record tiers, and the classical per-query and FFT costs.

% figure environment generated by tex_backend from the figure spec
% spec: fig:nf2 (status: banked)
\begin{figure*}[tb]
    \centering
    \includegraphics[width=\textwidth]{figures/NF2.pdf}
    \caption{Panel (a) shows the Gaussian acquisition aperture over the retained Gauss-Hermite nodes - scaled Hermite roots whose density is visibly non-uniform - with the Plancherel-Rotach retention threshold marked. Panel (b) shows the supported $(s, \tau)$ domain of one banked record with its three boundary causes drawn where they act: maximum evolution time, node retention, and the shot-noise floor. Both panels are schematics. Panel (c) plots the retained-node count against the law $\bar m \sim \|H\|\sqrt{\tau}$ using exact diagonalisation data from a single instance per system size (data source ED-N1r), with the analytical prediction drawn as a prediction and the fitted scaling drawn separately as a fit. Panel (c) uses a single instance per system size.}
    \label{fig:nf2}
\end{figure*}

% table environment generated for AO2 (canonised-run): the independent acquisition
% resources of one SSG campaign, collected plainly near the end of Sec IV. Wide cells
% use in-cell \parbox (not a p{..} preamble arg) so the brace-free {lll} spec parses
% to three columns and the long prose cells do not run overfull at body font size.
\begin{table*}[tb]
\caption{Independent acquisition resources for one Spectral Surface Generation campaign supporting $M$ post-acquisition spectral locations at inverse scale $\tau$. Here $\|H\|$ is the spectral-bandwidth bound, $\epsilon_w$ the aperture truncation tolerance, $K$ the shots per retained-overlap estimator and $C_{\tau,s}$ the reconstruction coefficient norm. Time points, evolution reach and record tiers are fixed at acquisition - the classical costs are per already-banked record. The table collects the governing relations already derived in the body and makes no new claim.}
\label{tab:resource}
\begin{ruledtabular}
\begin{tabular}{lll}
\textbf{Resource} & \textbf{Scaling or value} & \textbf{Governing relation} \\ \hline
\parbox[t]{0.22\textwidth}{\raggedright Retained time points $\bar m$\strut} & $\dfrac{t_{\max}}{\Delta t}\sim\dfrac{2}{\pi}\,\|H\|\sqrt{\tau\log(1/\epsilon_w)}$ & Eqs.~\eqref{res:aperture-resolution}, \eqref{res:bandwidth-grid} \\[6pt]
\parbox[t]{0.22\textwidth}{\raggedright Shots per retained estimator\strut} & $K$, additive error $\varepsilon_K=c\,K^{-1/2}$ & Eq.~\eqref{def:shot-budget-K} \\[6pt]
\parbox[t]{0.22\textwidth}{\raggedright Maximum evolution time\strut} & $t_{\max}=2\sqrt{\tau\log(1/\epsilon_w)}$, branch depth $t_{\max}/2$ & Eqs.~\eqref{res:aperture-resolution}, \eqref{res:half-depth-symmetric-branches} \\[6pt]
\parbox[t]{0.22\textwidth}{\raggedright Denominator record\strut} & \parbox[t]{0.40\textwidth}{\raggedright one-sided complex $A_j$, negative times by conjugation - the minimal bank\strut} & Eqs.~\eqref{res:one-sided-reconstruction}, \eqref{def:record-architecture-three-tier} \\[6pt]
\parbox[t]{0.22\textwidth}{\raggedright Branch records\strut} & \parbox[t]{0.40\textwidth}{\raggedright commuting tier one-sided; non-commuting $B_O$ two-sided, costed separately\strut} & Eq.~\eqref{def:record-architecture-three-tier} \\[6pt]
\parbox[t]{0.22\textwidth}{\raggedright Total shot budget\strut} & \parbox[t]{0.40\textwidth}{\raggedright $O(\bar m\,M\,K)$ over $M$ locations, with $K\gtrsim C_{\tau,s}^2/(\epsilon^2 Z_\tau^2)$ for an absolute read\strut} & Eq.~\eqref{res:shot-floor-visibility-wall} \\[6pt]
\parbox[t]{0.22\textwidth}{\raggedright Classical cost per query\strut} & $O(\bar m)$ per $(s,\tau)$ evaluation & Sec.~\ref{sec:N3} \\[6pt]
\parbox[t]{0.22\textwidth}{\raggedright Classical cost on an $s$ grid\strut} & \parbox[t]{0.40\textwidth}{\raggedright $O(\bar m M)$ direct, $O((\bar m+M)\log(\bar m+M))$ by FFT; $L$ scales cost $L$ transforms\strut} & Sec.~\ref{sec:N3} \\
\end{tabular}
\end{ruledtabular}
\end{table*}

% PLAN:N5
% NODES: goal:perturbation-assisted-itqde def:answer-key def:subset-sum-instance def:integer-weights def:subset-sum-function def:hamiltonian-h0 irf:h0-encodes-the-instance res:subset-sum-spectral-encoding def:degeneracy def:broad-support-state irf:subset-sum-as-spectral-detection def:degeneracy-profile irf:plus-state-overlap-sum irf:product-form-factorisation thm:hardy-littlewood-collapse irf:fourier-coefficient-of-g irf:circle-method-equivalence res:two-barriers-separated def:bounded-noncommuting-V irf:noncommutativity-breaks-product-form irf:single-link-suffices-to-break-chain def:cluster-projector-Ps def:isolation-gap-Gamma_S irf:weyl-stability-of-cluster irf:schur-complement-cluster-effective def:cluster-effective-hamiltonian def:internal-perturbation-block irf:internal-block-spectrum-governs-splitting def:cluster-substructure res:perturbative-rescue def:plateau-tilt irf:plateau-tilt-multiplicity-fingerprint def:multiplicity-fingerprint thm:conditional-visibility-lift def:null-channel irf:zz-first-order-fingerprint irf:free-fermion-trap-quadratic-v res:j-dial-fingerprint-scaling def:tilt-separation-ratio res:answer-key-certificate res:degeneracy-cartography-single-family res:operationally-dense-target-window irf:staircase-regime-broader-than-b-decision irf:step-displacement-log-ratio def:degeneracy-anchor irf:chained-anchor-degeneracy-profile res:integer-decoder def:integer-snap-and-trial-median def:decoder-gap-scales res:dial-feasibility-window def:erasure-rule def:operating-window
\section{The controlled subset-sum embodiment}\label{sec:N5}

% N5-p1
The controlled embodiment programs a subset-sum instance on integer weights $(a_1, \dots, a_n)$ with target $T$ - positive integer weights at the production scale $a_i = O(n^2)$, and for generic instances distinct weights with no planted structure - and reads every output of the method against an independently computable classical comparator that, for a given instance and known perturbation regime, sorts the measured surface into three classes of check,
\begin{widetext}
\[
Z_\tau(s)\ \text{checked}\
\begin{cases}
\text{exactly,} & \text{by dynamic programming},\\
\text{perturbatively,} & \text{through the modelled $V$-dressing},\\
\text{statistically,} & \text{by finite-shot validation}.
\end{cases}
\]
\end{widetext}
The family is the commuting classical-collapse control of the method: it supplies an exact reference at the production weight scale, a controlled degeneracy structure, and a graded route from the commuting limit into a genuinely noncommuting regime. The reconstruction and response identities on the supported $(s, \tau)$ domain, the shared-record covariance and the device's native execution of $H_0 + V$ are the machinery this instance now exercises, and the instance itself enters the protocol solely through the diagonal encoding.

% N5-p2
An ordered tuple $a = (a_1, \dots, a_n)$ of nonnegative integers identifies each bit string $x \in \{0,1\}^n$ with the subset $\{i : x_i = 1\}$, and the subset-sum function $s(x) = \sum_i a_i x_i$ takes values in $\{0, 1, \dots, \sum_i a_i\}$. The weights are carried by the diagonal Hamiltonian
\[
H_0 \;=\; \sum_{i=1}^{n} a_i\, \frac{I - Z_i}{2},
\]
each summand $(I - Z_i)/2$ being the single-qubit number operator - it returns $0$ on $|0\rangle_i$ and $1$ on $|1\rangle_i$ - so the instance lives in single-qubit drive amplitudes. The target $T$ enters through the probe shift $H_0 - T$. The instance therefore enters the protocol entirely through the Hamiltonian, and the initial state remains the maximally uncommitted broad-support reference, instance independent by construction: nothing measured later can be attributed to instance information smuggled through state preparation.

% N5-p3
The encoding surfaces in the spectrum. $H_0 |x\rangle = s(x)\, |x\rangle$ for every computational basis string, so the spectrum of $H_0$ is the multiset of subset sums of the weights, and the eigenspace at energy $S$ has dimension
\[
g(S) \;=\; \bigl|\{x \in \{0,1\}^n : s(x) = S\}\bigr|,
\]
the number of bit strings with subset sum $S$. The broad-support probe $|\psi_+\rangle = |+\rangle^{\otimes n} = 2^{-n/2} \sum_x |x\rangle$ places equal spectral weight $2^{-n}$ on every string, and under any initial state with broad computational support - the equal-weight probe is the operational default, not a requirement of the equivalence - the subset-sum decision problem is equivalent to detecting nonzero spectral weight at the target energy, $g(T) \ge 1$. Degeneracy is the meaning of this spectrum, and every readout that follows addresses it.

% N5-p4
An exact classical computation of $g$ on these instances costs no more than their production scale allows. The subset-sum decision problem is NP-complete in general, with best-known classical time scaling as $2^{0.3113\, n}$ at density one \cite{hgj2010} and best-known quantum time scaling as $2^{0.236\, n}$ \cite{bonnetain2020improved}, both up to polynomial factors, and heuristic differencing addresses the optimisation variant \cite{karmarkar1982differencing}. Those exponentials describe the general problem, not the programmed instances. At the production scale $a_i = O(n^2)$ the lattice of reachable sums is polynomially wide, so the Bellman counting recursion \cite{bellman1957dp} computes the exact count at every target in polynomial time (Fig.~\ref{fig:f5-dp-price}) - the price is itemised in Table~\ref{tab:dp-price}. The final row of the recursion is the degeneracy profile
\[
T \;\longmapsto\; g(T) \qquad \text{over occupied targets}\ T,
\]
absolute by construction, and this exact table is the validation reference for every reconstructed quantity in the section.

% N5-p5
At $V = 0$ the overlap record collapses exactly, and the derivation is two steps. From $|\psi_+\rangle = 2^{-n/2} \sum_x |x\rangle$ and the diagonal evolution $e^{-i t H_0} |x\rangle = e^{-i t s(x)} |x\rangle$, orthonormality of the computational basis gives $\langle\psi_+| e^{-i t H_0} |\psi_+\rangle = 2^{-n} \sum_x e^{-i t s(x)}$. The sum factorises over qubits, $\sum_x e^{-i t \sum_i a_i x_i} = \prod_i (1 + e^{-i t a_i})$, and with the target shift the record is
\[
\langle\psi_+|\, e^{-i t (H_0 - T)}\, |\psi_+\rangle \;=\; e^{i t T}\, 2^{-n} \prod_{i=1}^{n} \bigl(1 + e^{-i t a_i}\bigr),
\]
a product of single-qubit factors (Fig.~\ref{fig:nf3}). The single global normalisation is $2^{-n}$ - one factor of $1/2$ per qubit from the plus-state probe, equivalently the per-factor form $\prod_i (1 + e^{-i t a_i})/2$ - and never a single factor of $1/2$. The identity holds exactly at $V = 0$ with the plus-state probe. It is the baseline the perturbed protocol departs from, not a property of the perturbed record, and it is the exactly computable object every perturbed diagnostic in this section is read against.

% N5-p6
For integer weights the product expands as the generating Fourier series $\prod_i (1 + e^{-i t a_i}) = \sum_S g(S)\, e^{-i t S}$, and Fourier inversion at integer frequency gives
\[
g(S) \;=\; \frac{1}{2\pi} \int_{-\pi}^{\pi} dt\; e^{i t S} \prod_{i=1}^{n} \bigl(1 + e^{-i t a_i}\bigr).
\]
This integral is exactly a Hardy-Littlewood circle-method integral \cite{vaughan1997hardy}: the product is the generating exponential sum of $\{0,1\}^n$, and the coefficient at $S$ counts the representations of $S$ as a $\{0,1\}$-combination of the weights. The diagonal record therefore computes an exact Fourier coefficient of a finite exponential sum by real-time evolution, and the degeneracy profile is available from the $V = 0$ record in closed form (derived in full in App.~\ref{app:NA1}) - the exact reference at every target. The same fact read the other way: the factorised object is a product of single-qubit factors, classically free to evaluate, so any quantum content the full protocol carries must originate in features absent from it.

% N5-p7
The commuting picture's readout bottleneck separates into two independent quantities. Geometric isolation asks for a filter width that suppresses out-of-window weight by a factor $\eta$,
\[
\tau \;\ge\; \frac{\log(P_{\mathrm{out}}/P_{\mathrm{in}}) + \log(1/\eta)}{\Delta^2 - r^2}
\]
for intra-window radius $r$ and out-of-window gap $\Delta$ - polynomial exactly when $\Delta^2 - r^2 \ge 1/\mathrm{poly}(n)$ and $P_{\mathrm{out}}/P_{\mathrm{in}} \le \mathrm{poly}(n)$. Shot-floor visibility asks for a budget $K \ge c^2 \eta^{-2} P_{\mathrm{in}}^{-2}$ for a centred probe, with $c$ the retained-overlap coefficient norm, and in the commuting baseline $P_{\mathrm{in}} = g(T)/2^n$ for an exact target line, so a sparse target carries exponentially small visibility even at order-one integer spacing. A polynomial filter width that isolates the window does not by itself make the window visible above the sampling floor. The two barriers are independent and must be cleared separately - the perturbed diagnostics ahead change what is measured within a window and are never to be read as removing the visibility barrier.

% N5-p8
A bounded Hermitian perturbation noncommuting with the encoding breaks the collapse. We take $V$ local and two-body,
\[
\begin{gathered}
V \;=\; \sum_{a<b} V_{ab}, \qquad \|V\| < \infty, \\
[V, Z_i] \;\neq\; 0 \ \ \text{for some}\ i,
\end{gathered}
\]
with each term supported on two qubits. The factorisation used $e^{-i t s(x)} = \prod_i e^{-i t a_i x_i}$ - the exponential of a sum splitting as a product of exponentials of summands - and with noncommuting $V$ the exponential no longer splits, $e^{-i t (H_0 + V)} \neq e^{-i t H_0} e^{-i t V}$ in general. A single noncommuting matrix element suffices to fail the $n$-step factorisation chain, because the evolution mixes basis states across $s$-levels and the overlap acquires cross terms no single-qubit product reproduces. The break is generic rather than delicate - commutativity is not a robust assumption. But the break by itself is not yet a signal, and what it yields is quantified only after the cluster reduction.

% N5-p9
The perturbed spectrum is organised on two objects. The cluster projector
\[
P_S \;=\; \sum_{x \,:\, s(x) = S} |x\rangle\langle x|, \qquad Q_S \;=\; I - P_S,
\]
is the spectral projector of $H_0$ onto the subset-sum-$S$ eigenspace of dimension $g(S)$, and the projectors over distinct $S$ are mutually orthogonal idempotents summing to the identity. The isolation gap $\Gamma_S = \mathrm{dist}\bigl(S, \mathrm{spec}(Q_S H_0 Q_S)\bigr)$ is the spectral distance from $S$ to the nearest distinct subset-sum value, and for integer weights $\Gamma_S \ge 1$. The unit lower bound is a property of the integer encoding and not of the device, and it is the controlling parameter of every perturbative construction that follows.

% N5-p10
Whether the perturbation leaves this cluster structure intact is settled in two cases. Case (a), the crude global bound. Weyl's inequality \cite{bhatia1997matrix} confines every perturbed eigenvalue to within $\|V\|$ of its subset-sum value, and for the extensive local perturbation $V = J \sum_i V_i$ with per-term norm at most one and $m = O(n)$ bonds, subadditivity gives only $\|V\| \le mJ$, so disjointness of neighbouring discs demands $J = O(\Gamma_S/n)$ and the safe coupling shrinks with system size. Case (b), the quasi-local condition. Assume range-bounded local terms of coordination number $z$ and a differential diagonal shift $\delta_S = o(\Gamma_S)$ in the perturbative regime $J \ll \Gamma_S$. Then the block that actually couples the cluster to distinct values obeys $\|P_S V Q_S\| \le J \sqrt{z}\, \max_i \|V_i\|$, an $n$-independent quantity, and isolation holds whenever the effective gap $\Gamma_S - \delta_S$ exceeds twice this norm. Case (a) is sufficient and pessimistic rather than tight, because the extensive $O(nJ)$ displacement is dominated by the diagonal part of $V$, which shifts whole subset-sum sheets in common mode without mixing them. The operative condition is governed by the intensive ratio $J/\Gamma_S$ and the local connectivity, not by $n$, and case (b) holds under its stated assumption set.

% N5-p11
For an isolated cluster with $\Gamma_S > 4\|V\|$, the Riesz-projector and Schur-complement construction \cite{kato1995perturbation} reduces the full perturbed problem to a $g(S)$-dimensional block plus a bounded correction. The cluster effective Hamiltonian is
\begin{widetext}
\[
H_{\mathrm{eff}}^{(S)}(z) \;=\; S\, P_S + P_S V P_S + \Sigma_S(z), \qquad \Sigma_S(z) \;=\; P_S V Q_S\, \bigl(z - Q_S H Q_S\bigr)^{-1} Q_S V P_S,
\]
\end{widetext}
defined for $z$ in the cluster window $|z - S| < \Gamma_S/2$ and with $\|\Sigma_S(z)\| = O(\|V\|^2/\Gamma_S)$. The perturbed cluster eigenvalues read $E_{S,\alpha} = S + \nu_\alpha + O(\|V\|^2/\Gamma_S)$, where $\{\nu_\alpha\}$ are the eigenvalues of the internal block. The construction is valid in the stated analyticity window, the second-order self-energy is carried as a bounded correction rather than neglected, and this is the form in which every fingerprint statement is derived.

% N5-p12
The first-order physics lives in one Hermitian matrix, the internal perturbation block of dimension $g(S) \times g(S)$,
\begin{widetext}
\[
P_S V P_S \;=\; \bigl[\, \langle x | V | y \rangle \,\bigr]_{s(x) = s(y) = S}, \qquad \mathrm{spec}\bigl(P_S V P_S\bigr) \;=\; \{\nu_\alpha\}_{\alpha=1}^{g(S)},
\]
\end{widetext}
whose eigenvalues are the first-order intra-cluster splitting. The splitting pattern is determined by the matrix elements $\langle x|V|y\rangle$ for $x$ and $y$ inside the cluster - the cluster-restricted adjacency of $V$ on the solution set. Generic $V$ gives all distinct $\nu_\alpha$, while structured $V$ can preserve sub-degeneracies. The cluster substructure is the data $(g(T), \{p_\alpha\}, \{E_\alpha\})$ of the perturbed cluster - its degeneracy, the probe weights on the perturbed eigenvectors and their energies - equivalently the spectrum and eigenvector overlaps of the cluster effective Hamiltonian. What every response reads is this combinatorial geometry and not the multiplicity by itself: equal-multiplicity clusters can split differently, and structured perturbations can leave sub-degeneracies in place.

% N5-p13
A bounded local noncommuting $V$ with $\|V\| < \Gamma_S/4$ splits the $g(T)$-fold target eigenspace into sub-eigenvalues within the Weyl-stable isolation gap, and the within-cluster spread feeds a scalar ratio observable - the excess tilt at the target - measurable at polynomial shot budget at any filter width satisfying the stated feasibility conditions, because the ratio form cancels the common visibility factors. Under a validated forward model giving the spread as a known combinatorial function of the multiplicity and the weights, the observation inverts to $g(T)$. Two cost remarks position the perturbed family and carry no claim for the benchmark instances: reproducing the perturbed overlap family classically by matrix-product simulation meets volume-law entanglement growth at generic interacting $V$ \cite{calabrese2005evolution,eisert2010area,schollwock2011density,pal2010mbl}, and classical enumeration of the spread costs of order $g$ operations. At $g = 1$ the mechanism does not apply - a singleton has no internal spread - and dynamic programming already solves the sparse regime of the control polynomially, so the boundary is a scope statement rather than a limitation.

% N5-p14
On an isolated cluster plateau the staircase derivative is flat at $V = 0$ and acquires a tilt proportional to the within-cluster variance of the splitting,
\[
\frac{dH}{d\lambda} \;=\; -2\tau\, \sigma_\tau^2(\lambda),
\]
evaluated near the target. The forward model behind this observable is a functional of four inputs: the multiplicity $g(T)$, the within-subspace matrix elements of $V$ through $P_S V P_S$ and at next order $\Sigma_S$, the geometry of the solution set through the cluster-restricted adjacency, and the probe weights on the perturbed eigenvectors. The multiplicity fingerprint is the diagnostic family built from this response - demand the measured $J^2$ tilt exactly when the compressed perturbation has nonzero within-cluster variance, together with the $g = 1$ null - and interacting perturbations are validated routes to the response rather than logical guarantees. The universal reading that tilt size reports degeneracy is retired. Two clusters of equal multiplicity can respond differently, $P_S V P_S$ proportional to the identity gives zero internal variance at $g > 1$, and a generically vanishing block pushes the leading spread to second order. The response depends on the projected perturbation, the solution-set geometry and the probe weights, never on degeneracy alone.

% N5-p15
The conditional form of the response is
\begin{widetext}
\[
\sigma^2 \;=\; c\, J^2 \ \ \text{at leading order} \quad\Longleftrightarrow\quad \operatorname{Var}_{\psi_+}\!\bigl[\, P_S V P_S \,\bigr] \;\neq\; 0,
\]
\end{widetext}
the variance being the within-cluster variance of the compressed perturbation on the probed state at target $T$. The $g(T) = 1$ channel has no within-cluster variance - the invariant $\|P_S \psi_+\|^2 = g/2^n$ is trivially rotation-free in one dimension - so its tilt variance is identically zero by construction, and the singleton serves as the background-free null. Interacting $V_{XXZ}$ is the validated generic route to the response, and interaction alone is not a sufficient condition: symmetry, vanishing matrix elements or proportional-to-identity compression can silence a cluster under an interacting perturbation. The criterion is the compressed variance rather than the interaction label.

% N5-p16
The exact null is a statement about the isolated projected model, and a finite-$\tau$, finite-shot measurement around a $g = 1$ level sits on a small strictly positive empirical floor with five leading contributions. Neighbour leakage contributes $(P_{\mathrm{neigh}}/P_{\mathrm{in}})\, e^{-\tau d^2}$, with $d$ the distance to the nearest occupied energy. The second-order self-energy shift, of order $J^2/\Gamma_S$, displaces the plateau and feeds apparent variance at order $(J^2/\Gamma_S)^2$ through shifted neighbour tails. Out-of-cluster admixture at amplitude $O(J/\Gamma_S)$ contributes at order $(J/\Gamma_S)^2$ times the squared energy offsets. Finite-shot noise on the variance estimator carries the visibility factor $1/P^2$ with $P = g/2^n$ - exponential in $n$ at fixed budget on a truly sparse singleton, polynomial in the dense band. Reconstruction systematics from quadrature truncation and grid differencing enter at order the squared grid step. The measured floor $\sigma^2_{\mathrm{floor}}$ is the maximum of these five terms,
\begin{widetext}
\[
\sigma^2\bigl(g{=}1\bigr) \;\equiv\; 0 \ \ \text{(isolated projected model)}, \qquad \hat\sigma^2\bigl(g{=}1\bigr) \;=\; \sigma^2_{\mathrm{floor}} \;>\; 0 \ \ \text{(finite $\tau$, finite $K$)},
\]
\end{widetext}
and the observed invariance of the null at the level of a few parts per thousand confirms that the construction-level zero survives to that floor. Every exact-null statement is qualified the same way - exact in the isolated projected model, a small nonzero floor in the experiment.

% N5-p17
The two perturbation families of the benchmark instantiate the conditional on opposite sides. For $V_{XXZ}$ the diagonal $ZZ$ part moves no string and needs no resonance: it supplies a first-order within-cluster spread $\Delta_x = J \sum_i (2x_i - 1)(2x_{i+1} - 1)$, so the tilt scales as $J^2$ times the solution-set variance of the nearest-neighbour pair correlation and is identically zero at $g = 1$. For the $XX$ family the operator $X_i X_{i+1}$ flips an adjacent bit pair, changing the sum by $\pm(a_i + a_{i+1})$ on the $00 \leftrightarrow 11$ channel and by $\pm(a_i - a_{i+1})$ on the $01 \leftrightarrow 10$ channel. Positive weights never resonate on the sum channel, generic distinct weights never resonate on the difference channel, so $P_S V_{XX} P_S = 0$ exactly at first order and the splitting is pushed to the second-order self-energy, with tilt at order $J^4$. The $XX$ family is therefore a derived first-order selection-rule null - a property of the encoding and the perturbation, not an absence of signal, and a degenerate adjacent weight pair restores a first-order element through the local trade. Under Jordan-Wigner the family is quadratic and classically simulable by determinant methods at all $J$ \cite{valiant2002matchgates,terhal2002fermions,jozsa2008matchgates} - a statement about cost and not about spectral spread, since the cluster does split at second order - so the interacting choice $V_{XXZ}$ is the minimal hardware-native route to the first-order channel.

% N5-p18
Measured scaling follows the derived channel laws (Fig.~\ref{fig:nf4}). Exact in-window tilt variance at degenerate targets shows log-log slopes close to two for the $ZZ$ channel and close to four for the $XX$ channel against the derived guides, while the $g = 1$ channel is identically zero in the exact read. Under sampling noise the discrimination statistic is the tilt separation ratio
\[
\frac{\sigma^2(g \ge 2)}{\sigma^2(g = 1)},
\]
reading 2.1-2.5 at $K = 10^5$ for $J \le 0.05$, near 1.5 at $J = 0.1$, and gone by $J = 0.15$ as the splitting outgrows the filter window. The operational demand is a parametric family - the $J^2$ response on degenerate targets together with the null at $g = 1$ - rather than a single number. The fitted exponents are the as-plotted values for this data bank, and the discrimination holds across the sizes, couplings and shot budgets measured here.

% N5-p19
For known $V$ every value of the measured surface has an exact classical reference - dynamic programming for the $V = 0$ content plus the perturbation series for the $V$-dressing - so the residual between device and reference localises device error per energy window and per filter width. Decoherence rounds steps, miscalibration displaces plateaus, leakage adds phantom plateaus. The epistemic class is certification under trusted physics and never adversarial proof of quantumness: the reference assumes the physical model of the device rather than defending against an adversary. Within that class the reference-keyed residual is finer-grained than deployed volume-law benchmarks \cite{boixo2018xeb}, and key accuracy and classical reproduction cost trade against each other in $J$, so the mode operates inside the stated operating window.

% N5-p20
The channels assemble into one protocol statement. One retained overlap family at $V \neq 0$, measured in $\mathrm{poly}(n) \cdot M\sqrt{\tau_{\max}}$ circuits and shots, delivers the spectral surface of the programmed instance as one function-valued object $Z_\tau(s)$ with both filter parameters classical dials: the occupancy map, the absolute degeneracy profile chained from anchors, and the per-target multiplicity fingerprints. The perturbation adds the fingerprint channel and the certificate at zero cost to the density channel, because the total cluster weight $\|P_S \psi_+\|^2 = g/2^n$ is invariant under any intra-cluster rotation, and the $V = 0$ family is a product of cosines serving as the built-in calibration ruler. This is a protocol statement for the programmed control instance under a known perturbation family - the sensing inversion of an unknown $V$ is model-armed inference through the validated forward model, not assumption-free readout.

% N5-p21
The polynomial-shot extraction holds under the staircase-extraction regime, the conjunction of polynomially controlled geometric, quadrature and measurement axes. Visibility $P_{\mathrm{in}}^{-1} \le \mathrm{poly}(n)$ is not among the conditions, because cluster locations, widths and window emptiness extract from ratio observables riding the moment identities. The regime is broader than the dense-target decision setting - moderately sparse and density-one instances with a polynomially resolved geometric axis are in scope even where $g(T)/2^n$ is exponentially small. The regime changes which observables are in reach and does not remove the absolute-readout floor: ratio reads remain conditioned on the record clearing the additive shot floor.

% N5-p22
The elementary instrument of the cartography is a crossing read. Between consecutive occupied targets $T_a < T_b$ with spacing $\Delta$ and midpoint $M$, the centred staircase attains the midpoint value at a filter position displaced by
\[
u \;=\; \frac{\ln\bigl(g(T_b)/g(T_a)\bigr)}{2\tau\Delta},
\]
the balance point of the two flanking Gaussian bumps - a pure shape observable in which prefactors, overlap norm and the $2^{-n}$ all cancel. The formula is exact under five attached assumptions, each carrying a first-order correction through $\delta u = \epsilon/(2\tau\Delta)$. The two flanking clusters dominate the filtered weight. Internal cluster widths are negligible, with the leading bias set by the width difference. Probe weights are proportional to the wanted degeneracies - a common multiplicative weight error is invisible to the ratio, and only the differential error biases the read. Farther clusters are suppressed, which requires $\tau \ge O(\ln n)$ for polynomially varying $g$. Record error is controlled, propagating $\mathrm{Var}(u) = (\rho_a^2 + \rho_b^2)/(2\tau\Delta)^2$ with $\rho_c$ of order $K^{-1/2}/P_c$. The read is shot-optimal near $\tau\Delta^2 = 2$ and crosses empty targets through the $\Delta$ scaling.

% N5-p23
Subset sum supplies two exact anchors free,
\[
g(0) \;=\; 1, \qquad g\Bigl(\textstyle\sum_i a_i\Bigr) \;=\; 1,
\]
the empty subset and the full set, and chaining the crossing links $y_j = \ln g(T_{j+1}) - \ln g(T_j)$ from an anchor gives the absolute log-degeneracy profile. All links are reconstructed from the one banked record inside its supported $(s, \tau)$ domain, so their errors correlate exactly as the shared-record covariance of the acquisition section prescribes. Consecutive crossings share the cluster amplitude at their common target, the link covariance is tridiagonal, and the interior noise telescopes exactly, $\mathrm{Var}(\hat{l}_m) = \rho_0^2 + \rho_m^2$, so with the exact anchor the chained estimate at any step carries only endpoint shot noise. A common systematic breaks the cancellation at order $L^2 \sigma_{\mathrm{sys}}^2$. The two-anchor closure residual detects it without correcting it, and generalised least squares over both anchors spreads the closure mismatch by inverse variance in an $O(L)$ banded solve. The chained read is first an estimation problem and only then a discrete decoding problem.

% N5-p24
The profile readout is a decoding problem rather than pure estimation \cite{hua1990pencil,kay1993estimation}: the maximum degeneracy in the production sweep is two, so each step's truth is a log-ratio of small integers drawn from a tiny hypothesis class admissible from the current chained value. The decoder makes two aggregation moves,
\begin{widetext}
\[
y_j^{\mathrm{snap}} \;=\; \arg\min_{\rho \,\in\, \mathcal{H}_j} \bigl|\, \hat{y}_j - \ln\rho \,\bigr|, \qquad \hat{g}(T) \;=\; \mathrm{median}\bigl\{ \hat{g}^{(1)}(T), \dots, \hat{g}^{(8)}(T) \bigr\},
\]
\end{widetext}
the per-step integer snap onto the admissible ratio set $\mathcal{H}_j$ from the exact $g(0) = 1$ anchor, and an eight-trial per-target median. Together they lift continuous-read accuracy from 27-37\% to 56-66\% and then to 86-100\% at a total budget of $8 \times 10^5$ shots, and the infinite-shot exact read is 100\% on all tested families once the dial feasibility window and the erasure rule are applied. The tiny hypothesis class is a property of the controlled test with maximum degeneracy two - the decoder is scoped to the control, not offered as a general integer decoder - and the quoted rates are measured on the stated families and budgets.

% N5-p25
The decoder's accuracy figures presuppose a per-step filter setting, and the rule separates two spectral scales a single dial conflates. The pair gap $dT$ between the two compared clusters sets a visibility ceiling, because the mid-gap signal must clear the truncation floor, and the distance $d_{\mathrm{neigh}}$ to the nearest third cluster sets an isolation floor, because third-cluster bleed must be suppressed:
\begin{widetext}
\[
\tau \;\ge\; \frac{2}{d_{\mathrm{neigh}}^{2}} \ \ \text{(isolation floor)}, \qquad \tau \;\le\; \frac{\ln(1/x_{\mathrm{thr}})}{(dT/2)^{2}} \ \ \text{(visibility ceiling)}.
\]
\end{widetext}
The feasible window is $[\mathrm{floor}, \mathrm{ceiling}]$ and the prescribed $\tau$ lives inside it. When the window is empty the step is erased: the flagged maximum-likelihood default ratio of one is emitted and erased steps stay erased. An erased step is a flagged default and never a measurement, and erasure is preferred because it keeps the downstream chaining sound: a step with no feasible filter width would otherwise feed a certified-garbage read into the chain.

% N5-p26
One dial remains, the coupling itself. The certificate operates in the measured interval
\begin{widetext}
\begin{equation}\label{def:operating-window}
J \;\in\; \bigl[ J_{\mathrm{disc}},\, J_{\mathrm{esc}} \bigr], \qquad J_{\mathrm{disc}} \;=\; \frac{1}{4(n-1)}, \qquad J_{\mathrm{esc}} \;\approx\; 0.1, \qquad J^{*} \;\approx\; 0.05,
\end{equation}
\end{widetext}
with $J_{\mathrm{disc}}$ about 0.023-0.036 across the sizes probed. Below the lower edge the dynamics is effectively free and the perturbative response is not yet resolved. Above the upper edge the splitting outgrows the filter window and the $J^2$ family merges clusters. Both edges are classically computable dial settings fixed before any measurement - computed from the instance, not read from the device - and the window is measured for the controlled family. Every diagnostic of the embodiment, the exact reference and the decoder protocol operate inside this interval, and the numerical validation inherits them as the objects its runs exercise.

% figure environment generated by tex_backend from the figure spec
% spec: fig:f5-dp-price (status: banked)
\begin{figure}[tb]
    \centering
    \includegraphics[width=\linewidth]{figures/F5_dp_price.pdf}
    \caption{The Bellman dynamic-programming table $g_i(s)$ for the didactic instance $a = (1, \ldots, 6)$: each cell counts the subsets of the first $i$ weights that sum to $s$, computed exactly by the counting recursion, and the highlighted final row is the degeneracy profile $g(T)$. The table is computed inline and is deterministic. At the production weight scale $a_i = O(n^{2})$ the lattice of reachable sums is polynomially wide, so dynamic programming produces the exact degeneracy profile in polynomial time, and this profile is what every reconstructed read is checked against.}
    \label{fig:f5-dp-price}
\end{figure}

% figure environment generated by tex_backend from the figure spec
% spec: fig:nf3 (status: banked)
\begin{figure*}[tb]
    \centering
    \includegraphics[width=\textwidth]{figures/NF3.pdf}
    \caption{Panel (a) is a schematic of the exact factorisation at $V = 0$: the product formula for the overlap record of the diagonal subset-sum Hamiltonian $H_0$ is rendered as algebra with a compact per-qubit phase diagram, next to the $V \neq 0$ case in which a bounded noncommuting perturbation breaks the product form. Panel (b) shows the degeneracy profile $g(T)$ of the didactic instance $a = (1, \ldots, 6)$, computed inline as an exact spectral sum, with an inset showing a possible splitting pattern of a degenerate cluster under perturbation. The instance is didactic and chosen for legibility, and the splitting inset is one possible pattern rather than a prediction.}
    \label{fig:nf3}
\end{figure*}

% figure environment generated by tex_backend from the figure spec
% spec: fig:nf4 (status: banked)
\begin{figure*}[tb]
    \centering
    \includegraphics[width=\textwidth]{figures/NF4.pdf}
    \caption{Perturbative fingerprint response for the controlled subset-sum family. Panel (a) plots the tilt response of degenerate clusters against coupling $J$ at $\tau = 4$ from exact diagonalisation - the measured scaling fits give $J^{2.02}$ for the ZZ probe and $J^{4.00}$ for the XX probe, drawn with derived $J^{2}$ and $J^{4}$ guides. Panel (b) plots the separation of per-instance median responses between degenerate targets ($g \geq 2$) and unique targets ($g = 1$) against $J$ for the exact channel and for synthetic finite-shot channels at $K = 10^{3}$, $10^{4}$ and $10^{5}$. Panels (c) and (d) plot per-target response distributions at $J = 0.05$ as empirical cumulative curves across seeded instances at $n = 8$ and $n = 10$ - the exact channel, whose $g = 1$ null is identically zero by construction, and the reconstructed dial channel with its empirical $g = 1$ floor at each shot budget. Data: v6/N13proto\_fingerprint.csv, v8/N13b\_fingerprint\_distribution.csv and v8/N13b\_separation.csv, master seed 20260701. Across the tested budgets the figure discriminates degenerate from unique targets, with the fitted exponents tracking the derived $J^{2}$ and $J^{4}$ channel laws over the plotted range.}
    \label{fig:nf4}
\end{figure*}

% table environment generated by tex_backend from the table spec
% spec: tab:dp-price
% bite-zero-B: counterfactual table rebuild lifted here. The wide table's regime
% and method columns wrapped via \parbox in-cell (not a p{..} preamble arg): the
% column spec is brace-free {llll} so it parses to 4 columns matching the 4 body
% cells (the malformed flag was a p{0.32\textwidth} brace defeating the tabular-spec
% parser), while the \parbox keeps the long regime and method cells from running
% overfull at the full body font size.
\begin{table*}[t]
\caption{Classical price families for computing the degeneracy profile of a subset-sum instance at the production weight scale $a_i = O(n^{2})$. Rows list the exact dynamic-programming count, the SSG stored-record footprint, and the enumeration algorithms - meet-in-the-middle, Howgrave-Graham and Joux, and the quantum-walk variant. Entries are closed-form resource scalings from the cited analyses, with the record footprint taken from the retention law $\bar m \approx 3.7\,\|H\|\sqrt{\tau}$ measured in ED-N1r. The table places the benchmark price: lattice methods scale with the spectral width $\|H\|$, enumeration methods scale as $2^{cn}$ and matter only when the weights make the lattice exponentially wide, and at the production weight scale dynamic programming is exact and polynomial, with the SSG stored record riding the same lattice axis.}
\label{tab:dp-price}
\begin{ruledtabular}
\begin{tabular}{llll}
\textbf{Method} & \textbf{Price family} & \textbf{Leading-order scaling} & \textbf{Regime where the price binds} \\ \hline
\parbox[t]{0.2\textwidth}{\raggedright Dynamic programming (Bellman table)\strut} & lattice & $n\,\|H\|$ & \parbox[t]{0.32\textwidth}{\raggedright exact and polynomial at the production weight scale $a_i = O(n^2)$\strut} \\
\parbox[t]{0.2\textwidth}{\raggedright SSG stored record\strut} & lattice & $\bar m \approx 3.7\,\|H\|\sqrt{\tau}$ & \parbox[t]{0.32\textwidth}{\raggedright the same lattice axis, at the retention-law footprint\strut} \\
\parbox[t]{0.2\textwidth}{\raggedright Meet-in-the-middle\strut} & enumeration & $2^{n/2}$ & \parbox[t]{0.32\textwidth}{\raggedright weights exponentially wide in $n$\strut} \\
\parbox[t]{0.2\textwidth}{\raggedright Howgrave-Graham and Joux\strut} & enumeration & $2^{0.337n}$ & \parbox[t]{0.32\textwidth}{\raggedright weights exponentially wide in $n$\strut} \\
\parbox[t]{0.2\textwidth}{\raggedright Quantum walk\strut} & enumeration & $2^{0.216n}$ & \parbox[t]{0.32\textwidth}{\raggedright weights exponentially wide in $n$\strut} \\
\end{tabular}
\end{ruledtabular}
\end{table*}

% PLAN:N6
% NODES: goal:window-hunt res:answer-key-certificate res:integer-decoder def:integer-snap-and-trial-median def:erasure-rule def:decoder-gap-scales irf:window-entangling-but-localized irf:integrability-broken-by-disorder irf:n14-window-chi-flat-through-n20 irf:per-target-join-read res:window-verdict def:operating-window res:j-dial-fingerprint-scaling res:dial-feasibility-window
\section{Numerical validation}\label{sec:N6}

% N6-p1
The campaign answers one question: does a coupling window exist in which the reconstruction still agrees with the reference while the dynamics under $V_{\mathrm{XXZ}}$ is entangling, the convergence boundary set by local resonances $|a_i - a_j|$ rather than the global Weyl disc? It does. A nonempty operating window is measured below, and inside it reconstruction, the multiplicity fingerprint and the integer decoder all hold while the perturbation is entangling. Every reconstructed quantity is compared against its independently computable reference, the dynamic-programming lattice content dressed by the perturbation series. Each result is obtained by one of four methods - exact algebra evaluated numerically, exact diagonalisation at $n \le 12$, synthetic finite-shot sampling at stated shot budget, or TEBD at $n \le 20$. The banked data and the machine manifest resolving each data-source identifier are catalogued in App.~\ref{app:NA2}.

% N6-p2
Under synthetic finite-shot sampling over the benchmark instance set, data source SYN-N12c, decoder coverage is the $K$-independent geometric erasure count: 252 of 264 steps feasible. A step whose dial feasibility window is empty - the pair-gap visibility ceiling below the third-cluster isolation floor - is recorded as an erasure carrying the flagged ratio-1 maximum-likelihood default, excluded from conditional accuracy, and stays erased. Figure~\ref{fig:nf5} scores the feasible steps against the reference: the decoder snaps chained ratio estimates to the small-integer hypothesis class, aggregates by trial median, and should gain accuracy as shot noise falls below the gap scales. Single-trial and trial-median conditional accuracy are reported against shot budget $K$, full-profile recovery against $K$, and the ratio-1 imputation variant as its own clearly separated curve, with 95\% bootstrap intervals throughout. The claim is benchmark-specific - estimation then constrained decoding over a controlled hypothesis class, maximum degeneracy two - and reaches neither readability below the visibility floor, decoding beyond the tested class, nor performance outside the benchmark family.

% N6-p3
The dynamical half of the question is measured by exact diagonalisation at $n \in \{8, 10, 12\}$, data source ED-N11. Under $V_{\mathrm{XXZ}}$ the dynamics is entangling from the perturbative disc edge $J_{\mathrm{disc}} = 1/(4(n-1))$ of the operating window~\eqref{def:operating-window} upward, with required bond dimension $\chi_{\mathrm{req}}$ of 3 to 4 and saturated half-cut entanglement entropy $S_{\mathrm{vN}}$ of 0.5 to 1.3. Both quantities are flat in $n$ - $n = 8$ matches $n = 12$ to two decimals - which is the behaviour of localised cluster resonances rather than volume-law growth. The curves are measured per system size on the benchmark instance sequence and drawn as individual curves per size, never as a min-max band, and instance-to-instance bands are deferred.

% N6-p4
The operating model is a random-field XXZ chain - the diagonal disorder $a_i \sim n^{2}$ arises from the instance weights - so the Bethe-ansatz integrability of the clean XXZ chain is broken. The applicable frame is the localization physics of disordered spin chains, and that frame predicts both a flat-in-$n$ bond-dimension band and deepening localization as the disorder scale $W \sim n^{2}$ outruns a fixed coupling $J$. Under it the measured flatness is expected physics rather than a small-size artefact. The frame explains the band and generates a testable prediction.

% N6-p5
Sizes beyond $n = 12$ are inaccessible to exact diagonalisation, so we extend the measurement to $n = 20$ with TEBD in the Vidal gauge, the diagonal fields folded into the bond gates so that the Trotter commutator scales as $J^{2} a$ rather than $J a^{2}$. Validation precedes production: against the exact results at overlapping sizes the maximum entropy deviation is of order $10^{-6}$ to $10^{-7}$ (Fig.~\ref{fig:nf6b}), the required bond dimension agrees in 12 of 12 cells, and step-halving checks are flat. The production sweep over $n \in \{14, 16, 20\}$ and $J$ from 0.05 to 0.5 completes 24 of 24 cells with zero bond-dimension cap hits, data source TEBD-N14. Across $n = 8$ to 20 - a 2.5$\times$ span joining exact diagonalisation to TEBD - $\chi_{\mathrm{req}}$ stays in $\{2, 3, 4\}$, flat to falling, and the $S_{\mathrm{vN}}$ saturation is instance-dominated and $n$-independent. The flat-to-falling band is the deepening localization the disorder frame predicts. Past the exact-diagonalisation wall the TEBD results cover state properties only, and caps are reported as cap hits rather than truncated.

% N6-p6
The certificate's operating statement is read per degenerate target rather than globally in $J$: the fingerprint and the decoded profile are co-located on targets of degeneracy $g \ge 2$ by construction, so the join of decoding accuracy and fingerprint response is evaluated target by target. The per-target reading narrows the claim - a window statement quotes the targets on which both members of the join hold, and no global-in-$J$ guarantee is implied.

% N6-p7
The measured join is nonempty. The operating window runs from the perturbative disc edge $1/(4(n-1))$, numerically 0.023 to 0.036, up to $J \approx 0.1$, with operating point $J^{*} \approx 0.05$. Within it three measurements hold simultaneously. The dynamics is entangling and interacting - $\chi_{\mathrm{req}}$ of 3 to 4, $S_{\mathrm{vN}}$ of 0.5 to 1.3, from exact diagonalisation. The interacting-only tilt fingerprint clears the shot-noise floor under synthetic finite-shot sampling at production $K = 10^{5}$: separation 2.1 to 2.5$\times$ for $J \le 0.05$, 1.5$\times$ at $J = 0.1$, gone by roughly 0.15. The degeneracy profile decodes - 100\% in the unlimited-shot exact read, 82 to 89\% under synthetic finite-shot sampling at total budget $S_{\mathrm{tot}} = 8 \times 10^{5}$. Both window edges are classically computable dial settings - perturbative validity below, fingerprint escape above. Certificate mode operates throughout, with sensing mode the floor beneath it. The edges are approximate measured values on the benchmark family, the lower varying with $n$ through the disc-edge formula.

% N6-p8
Figure~\ref{fig:nf6} aligns three panels on the coupling axis $J$. The fingerprint under test is the interacting-only tilt channel, predicted to respond as $J^{2}$ on degenerate targets and escape above the window. Panel (a) shows the fingerprint separation from exact diagonalisation with synthetic finite-shot sampling overlaid, data source ED-N13. Panel (b) shows $S_{\mathrm{vN}}$ and $\chi_{\mathrm{req}}$ from exact diagonalisation at $n \in \{8, 10, 12\}$ as individual curves per system size, data source ED-N11. Panel (c) continues $\chi_{\mathrm{req}}$ to $n = 20$ with TEBD on an explicitly marked instance sequence, data source TEBD-N14, with $\chi_{\mathrm{req}}$ reported at the $10^{-4}$ discarded-weight convention under a per-gate truncation tolerance of $10^{-12}$. The alignment is the point - readability and simulation cost are statements about the same dial. Inside the readable window the fingerprint separates while the cost stays small and roughly $n$-independent, rising outside it. The figure records the operating window for the tested distinct-weight instances.

% figure environment generated by tex_backend from the figure spec
% spec: fig:nf5 (status: banked)
\begin{figure*}[tb]
    \centering
    \includegraphics[width=\textwidth]{figures/NF5.pdf}
    \caption{Decoder performance over the benchmark instance set under synthetic finite-shot sampling (data source SYN-N12c). Panel (a) shows coverage, the $K$-independent geometric erasure count, 252 of 264 steps feasible. Panel (b) shows conditional accuracy for single-trial and trial-median decoding. Panel (c) shows full-profile recovery against shot budget $K$, with 95\% bootstrap intervals throughout and the ratio-1 imputation variant drawn as its own clearly separated curve. Steps whose dial feasibility window is empty are recorded as erasures, excluded from conditional accuracy, and stay erased. Across the tested budgets the decoder recovers the degeneracy profile by estimation followed by constrained decoding over a controlled hypothesis class of maximum degeneracy two.}
    \label{fig:nf5}
\end{figure*}

% figure environment generated by tex_backend from the figure spec
% spec: fig:nf6 (status: banked)
\begin{figure*}[tb]
    \centering
    \includegraphics[width=\textwidth]{figures/NF6.pdf}
    \caption{Aligned panels share the coupling axis $J$. Panel (a) shows the fingerprint separation from exact diagonalisation with synthetic finite-shot sampling overlaid (data source ED-N13). Panel (b) shows the half-cut entanglement entropy $S_{\mathrm{vN}}$ and the required bond dimension $\chi_{\mathrm{req}}$ from exact diagonalisation at $n \in \{8, 10, 12\}$ (data source ED-N11), drawn as individual curves per system size. Panel (c) continues $\chi_{\mathrm{req}}$ to $n = 20$ with TEBD on an explicitly marked instance sequence (data source TEBD-N14), with $\chi_{\mathrm{req}}$ reported at the $10^{-4}$ discarded-weight convention under a per-gate truncation tolerance of $10^{-12}$. Across the tested distinct-weight instances the figure records the operating window: inside it the fingerprint separates while the simulation cost stays small and roughly $n$-independent, and the cost rises outside that range.}
    \label{fig:nf6}
\end{figure*}

% figure environment generated by tex_backend from the figure spec
% spec: fig:nf6b (status: banked)
\begin{figure*}[tb]
    \centering
    \includegraphics[width=\textwidth]{figures/NF6b.pdf}
    \caption{Matrix-product validation against exact diagonalisation in the time domain (data source ED/TEBD-N15b). Panel (a) overlays the half-cut entanglement entropy $S_{\mathrm{vN}}(t)$ from exact diagonalisation (line) and TEBD (circles) at $J = 0.2$ for $n \in \{8, 10, 12\}$, the two evolutions tracking one another across the record. Panel (b) plots the pointwise difference $|S_{\mathrm{vN}}^{\mathrm{ED}} - S_{\mathrm{vN}}^{\mathrm{TEBD}}|$ over all 24 production cells and two instances, staying below the marked maximum $2.7 \times 10^{-4}$. Across the overlapping sizes TEBD reproduces the exact entanglement dynamics, so the production sweep to $n = 20$ rests on a controlled matrix-product method.}
    \label{fig:nf6b}
\end{figure*}

% PLAN:N7
% NODES: res:sensing-grade-structural res:window-verdict irf:n14-window-chi-flat-through-n20 thm:readability-hardness-nogo irf:local-trades-absent-generic irf:denominators-grow-with-n irf:n14b-delocalization-only-far-above-ceiling def:localization-protection irf:stoquastic-cure-check irf:splitting-vs-clock def:planted-block thm:planted-escape res:planted-trade-family irf:xeb-epistemology-at-exponential-g fork:planted-hardness-formalisation
\section{Limits and discussion}\label{sec:N7}

% N7-p1
Within the validated operating window the tested generic distinct-weight instances are entangling but localized - the required bond dimension sits in $\{2, 3, 4\}$ and stays flat from $n = 8$ to $20$, and delocalization appears only far above the fingerprint ceiling. The benchmark therefore delivers a working spectral-sensing certificate on the generic family, its value standing on the fingerprint and decoder legs. The discussion separates three results the paper establishes at different strengths. The general acquisition-reconstruction architecture is exact and target-independent. The benchmark-specific diagnostics are measured over the stated window. The exploratory complexity extension of the planted family opens a question rather than settling one. The architecture is the widely applicable result. It asks only for a single coherent propagator expectation $G(t) = \langle\psi|e^{-iHt}|\psi\rangle$ - the quantity a universal quantum simulator produces \cite{lloyd1996universal, georgescu2014sim} and the real-time spectral estimators already sample \cite{somma2019eigenvalue, obrien2019qpe}. Any platform or protocol that can furnish that overlap therefore inherits the whole surface-and-response construction, with the mathematics carrying every later step on the banked record \cite{mccaul2025snacks, leamer2024emulation}. This section is where every global complexity boundary is stated in full.

% N7-p2
That flat band asks for an explanation, and we offer one as a finite-size compressibility proposition, hypotheses explicit. Take $H = H_{0} + JV$ on the generic distinct-weight family, with $V$ a bounded noncommuting perturbation of unit per-bond strength and $J$ inside the fingerprint window. Assume first the absence of local signed coincidences relevant to $V$, under which the off-resonance gap obeys $D_{\min} = \Omega(n)$ with probability $1 - \delta$ over the production ensemble. Assume second the local perturbation norm. Assume third a prethermal window $t^{*} \sim e^{c' n / J}$ covering the record horizon $t_{\max}(\tau)$~\eqref{res:aperture-resolution} \cite{abanin2017prethermal}. A single off-resonant rotation then makes the leading effective Hamiltonian diagonal, with
\begin{widetext}
\begin{equation*}
\begin{aligned}
&\lVert S^{(i)}\rVert \le \frac{J}{D_{\min}}\,, \qquad
\lVert R^{(i)}\rVert \le \frac{c\,J^{2}}{D_{\min}} = O\!\left(\frac{J^{2}}{n}\right), \qquad
\lVert R\rVert \sim \frac{n J^{2}}{D_{\min}} \sim J^{2}\,,\\
&\chi_{\varepsilon} = 2^{O(1)}\,\operatorname{poly}(1/\varepsilon) \quad \text{at fixed accuracy } \varepsilon\,, \ \text{independent of } n\,,
\end{aligned}
\end{equation*}
\end{widetext}
so the dressed residual is weak per bond while its global norm is a constant that does not vanish, and within the controlled window the physical state admits an efficient approximate matrix product state \cite{verstraete2008mps}. The proposition sits with the locality-bound and short-time literature, in which finite-window dynamics of local Hamiltonians is classically representable \cite{lieb1972finite, wild2023shorttime}. It is a finite-window compressibility statement about this construction, and it speaks to the tested range rather than to all classical representations or the worst case. The $e^{-S}$ undressing generally adds entanglement, so the physical state carries no exact bond-dimension-two guarantee.

% N7-p3
The quantitative engine of the proposition is combinatorial. On generic distinct-weight instances local trades - pairs of computational configurations differing on a bounded site set with equal instance energy - are absent, so every off-diagonal matrix element of $V$ is off-resonant by at least the local weight mismatch. The transition denominators are signed weight-sums over the move's window, $D_{xy} = \sum_{i} \epsilon_{i} a_{i}$. Distinctness alone secures $D_{xy} \neq 0$ for the two-site moves of $V_{\mathrm{XXZ}}$, where $D_{xy} \in \{\pm(a_{i} + a_{i+1}), \pm(a_{i} - a_{i+1})\}$, while a general $k$-local $V$ requires the explicit hypothesis that local signed coincidences relevant to $V$ are absent. Under that hypothesis, with the $a_{i}$ drawn i.i.d.\ from $\mathrm{Uniform}\{1, \dots, n^{2}\}$, a first-moment bound over the $O(n)$ nearest-neighbour differences gives $\Pr(D_{\min} \le \gamma) \le 2(n-1)\gamma/n^{2}$, so $D_{\min} \ge (\delta/2)\,n$ with probability $1 - \delta$ and the gap grows linearly in $n$. The bound is probabilistic over the production ensemble rather than a per-instance certainty.

% N7-p4
The other side of the coupling axis is measured rather than argued. The fingerprint ceiling is absolute - the filter window cannot pass couplings beyond roughly $J = 0.1$ to $0.2$ without merging unit-spaced clusters, and no $\tau$ setting rescues it - while delocalization requires $J$ at the local-resonance scale. Measured, $J = 1$ at ten times the ceiling gives bond dimensions $4/4/5$ across the sizes measured, a marginal onset, and $J = 3$ at thirty times the ceiling gives $6/6/11$ with the half-cut entropy $S_{\mathrm{vN}}$ rising from $0.96$ to $1.42$, so there the bond dimension grows with $n$. Entangling growth exists in this model, at couplings far above where the exact read can live, and the protocol's readable window does not reach those couplings.

% N7-p5
One mechanism accounts for both sides. The operating model is a random-field XXZ chain whose diagonal disorder comes from the instance weights,
\begin{equation*}
W \sim n^{2}\,, \qquad \frac{W}{J} \sim \frac{n^{2}}{J} \quad \text{at fixed in-window } J\,,
\end{equation*}
so at fixed coupling the disorder outruns the perturbation as $n$ grows. Localization deepens with size and in-window entanglement stays bounded - the flat bond dimension from $n = 8$ to $20$ - while raising $J$ towards the local-resonance scale closes the separation and produces the growth measured far above the ceiling. The picture is the established one for random-field XXZ chains: the numerically mapped localization of the random-field magnet \cite{znidaric2008mbl}, its phase transition \cite{pal2010mbl}, its description through local conservation laws \cite{serbyn2013liom} and its rigorous construction for spin chains \cite{imbrie2016mbl}. Because the disorder scale grows with $n$ by construction the family also sits near the strongly tilted systems studied as Stark localization, where fields that grow with system size localize without fixed-strength randomness \cite{schulz2019stark}, the localization a growing field induces in clean interacting chains \cite{vannieuwenburg2019stark} as observed without disorder \cite{morong2021stark}. Growing disorder is a regime the fixed-strength localization theorems do not directly cover, so we offer the mechanism as the explanation of the tested range, not as an asymptotic theorem about the family.

% N7-p6
The localization literature itself bounds what the mechanism can promise. The asymptotic status of many-body localization in one dimension is contested \cite{nandkishore2015mbl, abanin2019mbl}. Avalanches seeded by rare thermal regions can destabilise localization at large sizes and long times \cite{deroeck2017avalanche, luitz2017thermalize}, finite-size numerics are known to drift towards the thermal side as sizes grow \cite{morningstar2022avalanches}, and even deep in localized regimes entanglement grows logarithmically in time without bound \cite{bardarson2012unbounded}. Each of these failure modes concerns scales beyond the tested range, and none is observed within it. The caveats bound extrapolation beyond $n = 20$ and beyond the record horizon - they do not weaken the reported measurements.

% N7-p7
The family invites a check against the sign-problem literature, made here explicitly. The chain is one-dimensional, nearest-neighbour and bipartite - conjugating by $Z$ on the odd sublattice sends $XX + YY$ to $-(XX + YY)$ and leaves $ZZ$ and the diagonal $H_{0}$ invariant - so $H = H_{0} + J V_{\mathrm{XXZ}}$ is stoquastic up to a local basis change in the standard taxonomy \cite{bravyi2008stoquastic, klassen2019stoquastic}. Quantum Monte Carlo therefore faces no thermal sign problem in principle on this family \cite{troyer2005computational}, and because deciding whether a Hamiltonian can be cured of its sign problem is in general hard \cite{marvian2019cure, klassen2019hardness} we exhibit the transformation rather than assume it, in the spirit of methods that ease the sign problem where it cannot be removed \cite{hangleiter2020easing}. The surviving classical wall is observable access - the protocol quantities are filtered real-time-record objects, and extracting them from imaginary-time data is the ill-posed analytic-continuation problem, weaker than a sign problem and a statement about access to real-time observables rather than about the existence of an efficient equilibrium sampler.

% N7-p8
The generic bounds separate into two different scales. Readability constrains the instantaneous splitting scale $J\,B(L)$ - the staircase must stay inside the filter window at every size. Dynamics, and the entanglement it generates, depends on the integrated clock $Jt$. These are different quantities. A family that holds $J\,B(L)$ fixed while letting $Jt$ grow can therefore retain readability while its classical simulation cost grows, without contradicting the compressibility of the generic family. The pivot by itself produces no growth - the construction and the measurements that follow carry that burden.

% N7-p9
The escape family plants an equal-weight block of $L$ sites in the instance,
\begin{equation*}
a_{i} \equiv \bar{a} \ \text{ on the block}\,, \qquad H_{0}\big|_{\mathcal{M}} = E_{\mathcal{M}}\,, \qquad J = \frac{c}{L}\,,
\end{equation*}
where $\mathcal{M}$ is the block's degenerate manifold with common diagonal energy $E_{\mathcal{M}}$. Inside $\mathcal{M}$ the diagonal disorder cancels exactly, so $J$ acts there as a pure clock. The choice $J = c/L$ holds the worst in-window staircase spread bounded at every $n$ - the readability half - while entanglement accumulates on the integrated clock $Jt$, with measured bond dimension close to $2^{n/4}$. The construction follows the planted-instance tradition of hiding designed structure in an otherwise generic instance \cite{barthel2002planted}, and the degenerate manifold in which the clock runs is a designed analogue of the dynamically decoupled subspaces studied as scars and as Hilbert-space fragmentation \cite{turner2018scars, sala2020fragmentation}. The cancellation is exact only inside the planted manifold - outside it the generic disorder structure persists - and $J = c/L$ is a design choice of the family rather than a derived necessity.

% N7-p10
The theorem that results states the coexistence with both halves explicit,
\begin{widetext}
\begin{equation*}
J = \frac{c}{L}\,, \qquad J\,B(L) = 2JL = 2c \ \ \text{at every } n\,, \qquad Jt \ \ \text{growing on the protocol's own } \tau \text{ dial}\,.
\end{equation*}
\end{widetext}
With the block planted, in-manifold disorder cancels exactly and $J$ is a pure clock. The choice $J = c/L$ keeps the worst in-window spread bounded, so readability holds at every $n$, while the evolution time available to the protocol stretches polynomially on its own $\tau$ dial, so the integrated clock $Jt$ grows without breaching the window. The coexistence absent in the tested generic family - readability together with growing classical simulation cost - is realised by design. The theorem establishes that coexistence for the designed family over the analysed window, and it does not upgrade the generic verdict.

% N7-p11
Both halves of the trade carry numbers. Exact diagonalisation at $n = 12$, $L = 6$, $J = 0.033$ gives worst dressed in-window spread $0.413$, against both the $0.5$ threshold and the leading estimate $2JL = 0.4$, with the dressing shrinking as $O(1/n)$. TEBD at identical $(n, J, t)$ gives planted bond dimension $4/8/16/30$ at $n = 8/12/16/20$, close to $2^{n/4}$ with no cap hits, while the generic control stays pinned at bond dimension $2$ \cite{verstraete2008mps}. The clock cost is accounted through the acquisition ledger by the single substitution $t_{\max} = O(L)$ in the aperture reach~\eqref{res:aperture-resolution}: the supported filter-width domain widens to $\tau_{\max} = O(L^{2}/\log(1/\epsilon_{w}))$, the retained-node count grows as $\bar{m} = O(\lVert H\rVert L)$, the total shot budget as $S_{\mathrm{tot}} = O(\lVert H\rVert L)\,K$, and the per-window visibility floor stays flat in $L$ under a centred probe - the instantaneous spectral width is exchanged for a linearly longer clock at these derived acquisition costs, not obtained for free. Figure~\ref{fig:nf7} shows the trade: the worst in-window spread against $n$ with the $0.5$ threshold in panel (a), the required bond dimension for the planted family and the generic control with the $2^{n/4}$ guide in panel (b), the half-cut entanglement entropy against the accumulated clock $Jt$ in panel (c). The growth is consistent with exponential over the sizes measured. Panel (a) carries exact-diagonalisation points at $n = 8$ and $12$ only, exact diagonalisation being infeasible at $16$ and $20$, so readability at the larger sizes rests on the design bound rather than direct verification.

% N7-p12
At exponential degeneracy $g$ the per-step exact-integer read gives way to statistical certification against an exact prediction - the register of cross-entropy benchmarking \cite{boixo2018xeb}. The reference stays classically computable, because the planted block's in-manifold spectrum is exact by construction. The certificate changes epistemic register, not existence. That change bounds what any sampling-style comparison on the planted family could certify, and is part of the family's cost.

% N7-p13
The architecture the paper establishes is a reusable spectral-acquisition record and the family of reconstructions it supports, validated on the controlled benchmark across its measured operating window. Separately, the planted family opens a complexity question: whether the overlap record $A(t)$ on planted instances is formally hard to compute classically - the one question between the demonstrated family and a quantum-advantage-relevant claim. A single planted block is integrable-adjacent, its in-manifold dynamics close to a free precession, so the natural route is multi-block or junction constructions that harden the manifold, and we leave that formalisation as a named open problem. The best classical subset-sum algorithms set the benchmark any such claim would face \cite{bonnetain2020improved}. The observed bond-dimension growth is the cost of one simulation method over the analysed range, and whether it reflects a genuine classical lower bound is exactly what that open problem asks.

% figure environment generated by tex_backend from the figure spec
% spec: fig:nf7 (status: banked)
\begin{figure*}[tb]
    \centering
    \includegraphics[width=\textwidth]{figures/NF7.pdf}
    \caption{Panel (a) plots the worst in-window staircase spread against system size $n$ for the planted family, with the 0.5 readability threshold marked - exact-diagonalisation points at $n = 8$ and 12, and $n = 16$ and 20 explicitly marked as infeasible for exact diagonalisation. Panel (b) plots the required bond dimension $\chi_{\mathrm{req}}$ against $n$ for the planted family and a generic control from TEBD, with the $2^{n/4}$ line drawn and labelled as a derived guide. Panel (c) plots the half-cut entanglement entropy against the accumulated clock $Jt$ from TEBD. All matrix product state data use the linear chain ordering with site index 0 to $n - 1$ and block sites first, a per-gate truncation tolerance of $10^{-12}$, and a bond-dimension cap of 128 (data source ED/TEBD-N15b). Across $n = 8$ to $20$ the planted family realises the trade by design: readability holds while the required bond dimension grows, consistent with exponential over that range. The $2^{n/4}$ line is a guide, not an established law.}
    \label{fig:nf7}
\end{figure*}

% PLAN:N8
% NODES:
\section{Conclusion}\label{sec:N8}

% N8-p1
Spectral Surface Generation changes the unit of quantum work. Before it, a spectral protocol is organised around a requested readout and each new question is a new experiment. After it, one acquisition campaign produces a persistent time-domain object - the banked overlap record $\hat{G}(t_k)$ sampled at the Gauss-Hermite grid times - whose location-scale representation and responses are selected afterward. A classical layer reconstructs the Gaussian-filtered spectral surface $\hat{Z}_\tau(s)$, the generating surface and the response channels - filtered mean and variance, higher cumulants, within-window ratios - at post-acquisition $(s,\tau)$, each read a coefficient update on the stored record with no further quantum evolution. The reconstruction identity is an exact continuous transform, and the $\tau$ direction obeys a forward-stable heat-kernel scale space~\cite{witkin1983scale,koenderink1984structure,lindeberg1994scale}.

% N8-p2
One acquisition campaign supports reconstruction of $\hat{Z}_\tau(s)$ and its response channels over a finite-record envelope - the $(s,\tau)$ domain bounded by time reach from the Gaussian aperture, node retention from the Gauss-Hermite gate and the shot-noise floor. Shot noise propagates through the linear reconstruction with a derived covariance for shared-record ratio estimators, whose gain is a constant factor. The envelope is what one campaign actually acquires, and its outer reach is set by the visibility wall: an absolute read of exponentially small spectral weight costs exponentially many shots, so every readout is conditioned on its window clearing the shot floor.

% N8-p3
On the controlled subset-sum family every reconstructed quantity has an independently computable reference: at $V = 0$ the plus-state overlap factorises exactly and the degeneracy profile $g(T)$ is an exact Fourier coefficient of the record. The reconstruction agreed with that reference. The multiplicity fingerprint responded at order $J^2$ within its stated forward model - it reads the compressed perturbation's within-cluster variance, not $g$ itself - with the $g = 1$ null channel at its finite empirical floor. The integer decoder reported coverage and conditional accuracy under its erasure rule inside a nonempty measured operating window. The validation ran by exact diagonalisation at $n \le 12$, TEBD at $n \le 20$, and synthetic finite-shot sampling at stated budgets $K$.

% N8-p4
The planted family leaves open a single complexity question: whether its overlap records form a classically hard instance family, formalised through multi-block or junction constructions that harden the manifold while the banked record keeps supplying readable diagnostics. Beyond it the natural directions are broader kernels and observables, hardware acquisition, and harder embodiments. The capability the paper delivers is already concrete, and its reach is set by a single requirement: the method asks only for the expectation of a propagator, $G(t) = \langle\psi|e^{-iHt}|\psi\rangle$ - the Loschmidt amplitude an echo or interferometric protocol measures \cite{gorin2006loschmidt} and a controlled quantum simulator supplies \cite{georgescu2014sim}. Wherever that overlap can be acquired, a spectral experiment run once becomes a reusable classical record from which a whole supported family of spectral responses is read afterward, across spectral location and every coarser scale the acquisition reach allows. The mathematics does the rest.


\appendix
% PLAN:NA1
% NODES: def:hamiltonian def:computational-basis def:pauli-z def:plus-state def:imaginary-time-parameter def:filtered-partition ds:filter-in-trace-form ds:spectral-trace-collapse def:spectral-weights def:target-index-set def:in-out-spectral-weights ds:split-Z-by-window def:window-radius-and-gap ds:bound-Zout-above ds:bound-Zin-below def:staircase thm:hardy-littlewood-collapse ds:pauli-z-eigenequation ds:h0-acts-as-sum def:hamiltonian-h0 def:integer-weights def:subset-sum-function ds:exp-of-diagonal-h0 ds:expand-plus-tensor def:broad-support-state ds:sum-over-comp-basis ds:exponent-as-product ds:sum-product-interchange ds:evaluate-binary-inner-sum ds:product-as-fourier-series def:degeneracy thm:staircase-cumulant-identities ds:partial-lambda-log-Z def:filtered-mean def:centered-readout irf:retained-overlap-budget ds:gaussian-via-quadrature def:filter-kernel ds:lambda-shift-commutes def:hamiltonian-shift def:probe-target-shift ds:invert-for-tau-eta ds:quadrature-tail-decay def:quadrature-truncation-budget ds:invert-for-K-eta def:shot-budget-K def:cluster-effective-hamiltonian def:multiplicity-fingerprint ds:riesz-projector-cluster def:cluster-projector-Ps def:isolation-gap-Gamma_S def:bounded-noncommuting-V ds:block-diagonalise-via-similarity ds:eigenvector-as-superposition-on-cluster irf:internal-block-spectrum-governs-splitting ds:overlap-with-plus-state-coherent-sum def:perturbed-spectral-weights irf:target-cluster-second-order-spread res:perturbative-rescue
\section{Derivations}\label{app:NA1}

\subsection*{Primitive objects}

% NA1-p1
Every derivation below works from the same ground object - a Hermitian operator on the $2^n$-dimensional Hilbert space of $n$ qubits,
\begin{equation*}
H = H^\dagger = \sum_j E_j \, |E_j\rangle\langle E_j| ,
\end{equation*}
whose eigendecomposition over real eigenvalues $E_j$ and an orthonormal eigenbasis $\{|E_j\rangle\}$ is supplied by the spectral theorem. Every manipulation that follows either inserts this eigenbasis or diagonalises in it, so we put the decomposition on record before anything acts on it.

% NA1-p2
The concrete register basis is the computational basis - the $2^n$ product states
\begin{equation*}
|x\rangle = |x_1\rangle \otimes \cdots \otimes |x_n\rangle , \qquad x \in \{0,1\}^n ,
\end{equation*}
built from single-qubit $|0\rangle$ and $|1\rangle$. It is the joint eigenbasis of every $Z_i$ and the basis in which the subset-sum Hamiltonian is diagonal. The staircase and cluster derivations below manipulate amplitudes string by string, so the basis is fixed before any operator acts on it.

% NA1-p3
On a single qubit the Pauli operator is
\begin{equation*}
Z = \mathrm{diag}(+1, -1) , \qquad Z|0\rangle = +|0\rangle , \qquad Z|1\rangle = -|1\rangle ,
\end{equation*}
and $Z_i$ denotes $Z$ acting on qubit $i$ with the identity on the remaining qubits. The sign of $Z_i$ on $|x\rangle$ therefore encodes the bit $x_i$ - the convention from which the staircase derivation builds the subset-sum Hamiltonian.

% NA1-p4
The single-qubit state
\begin{equation*}
|+\rangle = \frac{|0\rangle + |1\rangle}{\sqrt{2}}
\end{equation*}
is the $+1$ eigenstate of Pauli $X$ and carries equal amplitude on both computational basis states. Its $n$-fold tensor power is the broad-support state the protocol prepares, and the overlap computations of the staircase and cluster derivations both expand it in the fixed basis.

% NA1-p5
The last primitive is the filter parameter,
\begin{equation*}
\tau > 0 , \qquad \Delta E_\tau \sim \tau^{-1/2} ,
\end{equation*}
a positive real number setting the Gaussian filter width. Larger $\tau$ sharpens spectral resolution, and it also grows the prefactor $\|H\|^{2\bar m}$ in the retained-overlap budget, so $\tau$ is the resource the truncation derivation inverts. Every filtered quantity from the collapse derivation onwards carries it without restatement.

\subsection*{Exact collapse of the filtered partition function and the window bounds}

% NA1-p6
The filtered partition function of the body,
\begin{equation*}
Z_\tau(\lambda) = \langle \psi | \, e^{-\tau (H - T + \lambda)^2} \, | \psi \rangle ,
\end{equation*}
collapses exactly onto the spectrum as $Z_\tau(\lambda) = \sum_j p_j \, e^{-\tau \xi_j^2}$ with $\xi_j$ the shifted energies defined below, and splitting the collapsed sum by the target window proves the two one-sided weight bounds. The body's filtered-partition display draws on the collapsed form, and the $\tau$ and $K$ inversions below draw on the two bounds.

% NA1-p7
With $\rho = |\psi\rangle\langle\psi|$ the rank-one density matrix of the initial state, the filtered expectation is a trace,
\begin{equation*}
Z_\tau(\lambda) = \langle \psi | \, e^{-\tau (H - T + \lambda)^2} \, | \psi \rangle
= \mathrm{Tr}\!\left[ e^{-\tau (H - T + \lambda)^2} \, \rho \right] .
\end{equation*}
The only manipulation is recognising the expectation as a trace against a rank-one projector, and trace cyclicity is automatic for this $\rho$. The trace form is the step that lets the resolution of identity act.

% NA1-p8
Inserting the resolution of identity $I = \sum_j |E_j\rangle\langle E_j|$ collapses the trace onto the spectrum,
\begin{widetext}
\begin{align*}
\mathrm{Tr}\!\left[ e^{-\tau (H - T + \lambda)^2} \, \rho \right]
&= \sum_j e^{-\tau (E_j - T + \lambda)^2} \, |\langle E_j | \psi \rangle|^2 \\
&= \sum_j p_j \, e^{-\tau (E_j - T + \lambda)^2} , \qquad p_j = |\langle E_j | \psi \rangle|^2 ,
\end{align*}
\end{widetext}
with $p_j$ the spectral weights of the initial state - real, non-negative and normalised, $\sum_j p_j = 1$. The filtered partition function therefore reads as a weighted Gaussian density of states sampled at the probe. The collapse is exact for every $\tau$ and every probe value - no approximation enters at this step. Every filtered observable in the paper is a functional of this sum.

% NA1-p9
Window membership is decided by the shifted energy $\xi_j(\lambda) = E_j - T + \lambda$ that the collapsed Gaussian factors carry. The target index set of radius $r$ is
\begin{equation*}
C = C(T, \lambda) = \left\{ j : |\xi_j(\lambda)| \le r \right\} , \qquad \xi_j(\lambda) = E_j - T + \lambda ,
\end{equation*}
the eigenvalue indices inside the target window centred at $T - \lambda$. The split of the collapsed sum needs this sharp membership rule before the partition is written.

% NA1-p10
The two bounds compare aggregates of the spectral weights split by membership in $C$,
\begin{equation*}
P_{\mathrm{in}} = \sum_{j \in C} p_j , \qquad
P_{\mathrm{out}} = \sum_{k \notin C} p_k = 1 - P_{\mathrm{in}} .
\end{equation*}
$P_{\mathrm{in}}$ is the spectral weight inside the target window and $P_{\mathrm{out}}$ the weight outside, and the pair splits the unit-normalised spectral weight exactly by window membership.

% NA1-p11
Partitioning the collapsed sum over the index set gives the exact two-term split
\begin{equation*}
Z_\tau(\lambda) = \sum_j p_j \, e^{-\tau \xi_j^2}
= \underbrace{\sum_{j \in C} p_j \, e^{-\tau \xi_j^2}}_{Z_\tau^{\mathrm{in}}}
+ \underbrace{\sum_{k \notin C} p_k \, e^{-\tau \xi_k^2}}_{Z_\tau^{\mathrm{out}}} .
\end{equation*}
Both parts are non-negative and the partition is exact - the only manipulation is a partition of the index set. The two sides are parametrised by the window radius $r \in [0, \Delta)$, with $|\xi_j| \le r$ inside, and the gap $\Delta$, with $|\xi_k| \ge \Delta$ outside. This is the geometry every later inequality refines.

% NA1-p12
The Gaussian is monotone decreasing in $|\xi|$, and that single fact yields both bounds. For $k \notin C$ the gap gives $e^{-\tau \xi_k^2} \le e^{-\tau \Delta^2}$, so
\begin{equation*}
Z_\tau^{\mathrm{out}} \le P_{\mathrm{out}} \, e^{-\tau \Delta^2} ,
\end{equation*}
and for $j \in C$ the radius gives $e^{-\tau \xi_j^2} \ge e^{-\tau r^2}$, so
\begin{equation*}
Z_\tau^{\mathrm{in}} \ge P_{\mathrm{in}} \, e^{-\tau r^2} .
\end{equation*}
The out-of-window contribution decays at the gap rate, and the in-window contribution retains at least the windowed weight times the worst-case intra-window suppression. Both bounds are one-sided and exact under the window hypotheses $|\xi_j| \le r$ inside and $|\xi_k| \ge \Delta$ outside - no distributional assumption on the spectrum enters. These are the in-window and out-of-window weight statements the body cites, and the resource inversions below draw on them.

\subsection*{Closed form of the free overlap as the subset-sum generating series}

% NA1-p13
The diagonal Hamiltonian stores its free evolution in the record in closed form. The body identity
\begin{widetext}
\begin{equation*}
\langle \psi_+ | \, e^{-i t H_0} \, | \psi_+ \rangle
= 2^{-n} \prod_{i=1}^{n} \left( 1 + e^{-i t a_i} \right)
= 2^{-n} \sum_{S \ge 0} g(S) \, e^{-i t S}
\end{equation*}
\end{widetext}
is derived here. The staircase definition of the record section and the Hardy-Littlewood collapse theorem both draw on the series, so each algebraic step below is read against the product-to-series statement it serves.

% NA1-p14
The algebra opens with the Pauli-$Z$ eigenequations fixed among the primitives. From $Z|0\rangle = +|0\rangle$ and $Z|1\rangle = -|1\rangle$,
\begin{align*}
Z_i \, |x\rangle &= (1 - 2 x_i) \, |x\rangle , \\
\frac{I - Z_i}{2} \, |x\rangle &= \frac{1 - (1 - 2 x_i)}{2} \, |x\rangle = x_i \, |x\rangle ,
\end{align*}
so each $(I - Z_i)/2$ acts as the per-qubit number operator, reading the bit $x_i$ off the string it meets. This identity converts every later operator statement into arithmetic on bits, and the rest of the walk is built on it.

% NA1-p15
Fix an ordered tuple $a = (a_1, \dots, a_n)$ of non-negative integers - the subset-sum instance - and let $s(x) = \sum_{i=1}^{n} a_i x_i$ denote the subset sum of the string $x$. Applying the diagonal Hamiltonian $H_0 = \sum_{i=1}^{n} a_i (I - Z_i)/2$ to a basis string and using linearity,
\begin{align*}
H_0 \, |x\rangle &= \sum_{i=1}^{n} a_i \, \frac{I - Z_i}{2} \, |x\rangle
= \left( \sum_{i=1}^{n} a_i x_i \right) |x\rangle \\
&= s(x) \, |x\rangle .
\end{align*}
Every computational basis string is an eigenstate of $H_0$ with eigenvalue equal to its subset sum under the weights $a$. This eigenvalue identity connects the quantum register to the combinatorial instance - the rest of the derivation only reorganises phases built on it.

% NA1-p16
Exponentiating turns the eigenvalues into phases. The eigenvalue identity gives the spectral decomposition $H_0 = \sum_{y} s(y) \, |y\rangle\langle y|$, whose exponential is
\begin{align*}
e^{-i t H_0} &= \sum_{y \in \{0,1\}^n} e^{-i t s(y)} \, |y\rangle\langle y| , \\
e^{-i t H_0} \, |x\rangle &= e^{-i t s(x)} \, |x\rangle ,
\end{align*}
the second line using only computational-basis orthonormality, $\langle y | x \rangle = \delta_{xy}$. Free evolution rotates each string by its subset-sum phase and transfers no amplitude between strings. This phase-only action is what makes the free overlap computable in closed form, so it is displayed before the overlap is assembled.

% NA1-p17
The second ingredient of the overlap is the initial state. Expanding the $n$-fold tensor power of the plus state in the fixed basis gives the uniform superposition
\begin{equation*}
|\psi_+\rangle = \left( \frac{|0\rangle + |1\rangle}{\sqrt{2}} \right)^{\otimes n}
= 2^{-n/2} \sum_{x \in \{0,1\}^n} |x\rangle ,
\end{equation*}
in which every basis string carries the same amplitude $2^{-n/2}$. This is the broad-support property of the initial state, displayed so the coherent sums that follow have explicit coefficients.

% NA1-p18
Sandwiching the diagonal evolution between two copies of the expanded state, factorising the phase and interchanging sum with product form one tight chain,
\begin{align*}
\langle \psi_+ | \, e^{-i t H_0} \, | \psi_+ \rangle
&= 2^{-n} \sum_{x \in \{0,1\}^n} e^{-i t s(x)} , \\
e^{-i t s(x)} &= e^{-i t \sum_i a_i x_i} = \prod_{i=1}^{n} e^{-i t a_i x_i} , \\
\sum_{x \in \{0,1\}^n} \, \prod_{i=1}^{n} e^{-i t a_i x_i}
&= \prod_{i=1}^{n} \, \sum_{x_i \in \{0,1\}} e^{-i t a_i x_i} .
\end{align*}
The first line uses computational-basis orthonormality to remove every off-diagonal term of the double sum. The second substitutes $s(x)$ and splits the exponential of a sum into a product of exponentials - legitimate because the scalar summands $a_i x_i$ commute. The third is the distributive law - the sum over strings factorises over independent coordinates because each summand factors into per-coordinate terms. The overlap is now a product of independent binary sums.

% NA1-p19
Each binary sum evaluates in one step, and collecting the product by subset-sum value closes the derivation,
\begin{align*}
\sum_{x_i \in \{0,1\}} e^{-i t a_i x_i} &= 1 + e^{-i t a_i} , \\
\langle \psi_+ | \, e^{-i t (H_0 - T)} \, | \psi_+ \rangle
&= e^{i t T} \, 2^{-n} \prod_{i=1}^{n} \left( 1 + e^{-i t a_i} \right) , \\
\prod_{i=1}^{n} \left( 1 + e^{-i t a_i} \right)
&= \sum_{S \ge 0} g(S) \, e^{-i t S} .
\end{align*}
The second line attaches the target prefactor $e^{i t T}$, which belongs to the recorded spectral signal rather than to the product identity. The third expands the product term by term - each summand corresponds to one choice of $x_i \in \{0,1\}$ per factor and contributes the phase $e^{-i t s(x)}$ - and regroups the $2^n$ terms by the value $s(x) = S$, so the coefficient at $S$ is the degeneracy $g(S) = |\{ x : s(x) = S \}|$. The result is a finite Fourier series in $t$, since $S$ ranges over the integers from $0$ to $\sum_i a_i$, with the subset-sum counts as its coefficients - the generating-function structure the Hardy-Littlewood collapse draws on \cite{vaughan1997hardy}.

\subsection*{The log-derivative readout identity}

% NA1-p20
The first staircase-cumulant identity of the body,
\begin{equation*}
\partial_\lambda \log Z_\tau(\lambda) = -2\tau \, m_\tau(\lambda) .
\end{equation*}
follows from one displayed derivative connecting the measurable partition function to the centred readout. The staircase-cumulant identities theorem uses this identity as its first entry, and the centred readout uses it as the zero condition that centres the filter on the target cluster.

% NA1-p21
Differentiating the collapsed spectral form of $\log Z_\tau(\lambda)$ in the probe, by the chain rule,
\begin{align*}
\partial_\lambda \log Z_\tau(\lambda)
&= -2\tau \, \frac{\sum_j p_j \, (E_j - T + \lambda) \, e^{-\tau (E_j - T + \lambda)^2}}{Z_\tau(\lambda)} \\
&= -2\tau \left( \mu_\tau(\lambda) - T + \lambda \right) \\
&= -2\tau \, m_\tau(\lambda) ,
\end{align*}
where $\mu_\tau(\lambda) = \sum_j p_j E_j \, e^{-\tau (E_j - T + \lambda)^2} / Z_\tau(\lambda)$ is the filtered mean and $m_\tau(\lambda) = \mu_\tau(\lambda) - T + \lambda$ the centred readout. The derivative is displayed in full because the body quotes only its endpoints - the intermediate line, a filtered average of the shifted energies, is where the identification with the filtered mean happens. The identity holds for every $\tau > 0$ and every probe value, and $m_\tau = 0$ is the centring condition on the filter rather than a stationarity assumption.

\subsection*{The finite record and the resource inversions}

% NA1-p22
The retained-overlap budget of the record section,
\begin{equation*}
\varepsilon \le \sum_{k > \bar m} w_k f_k , \qquad
\bar m = \frac{t_{\max}}{\Delta t} = \frac{\sqrt{8}}{\pi} \, \|H\| \sqrt{\tau \ln(1/\varepsilon_w)} ,
\end{equation*}
is derived here together with its companion inversions. The Gaussian filter acquires a finite unitary-overlap representation, the probe shift factors out of the record as a scalar phase, and demanding relative error $\eta$ yields the lower bounds on the filter strength $\tau$, the retained-pair budget $\bar m$ and the shot budget $K$ that the record and measurement sections use. The three inversions form one derivation because each converts the same accuracy demand into a lower bound on one resource.

% NA1-p23
The acquisition step behind the retained record converts the imaginary-time filter into real-time evolutions the hardware can run. The filter has the inverse Fourier representation
\begin{widetext}
\begin{equation*}
e^{-\tau (H - T + \lambda)^2} = \int dt \, K_\tau(t) \, e^{-i t (H - T + \lambda)} , \qquad
K_\tau(t) = \frac{1}{\sqrt{4 \pi \tau}} \, e^{-t^2 / 4\tau} ,
\end{equation*}
\end{widetext}
with a Gaussian kernel, and Gaussian quadrature at nodes $\{t_j\}_{j \in J}$ with weights $\{w_j\}$ truncates the integral to the finite unitary sum
\begin{equation*}
e^{-\tau (H - T + \lambda)^2} \approx \sum_{j \in J} w_j \, e^{-i t_j (H - T + \lambda)} ,
\end{equation*}
one recorded evolution per retained node. The representation itself is exact before truncation - truncation introduces the tail error whose decay and inversion later steps derive.

% NA1-p24
With $H_T = H - T I$ the target-shifted Hamiltonian, the truncated evolutions run under $H_T^{(\lambda)} = H_T + \lambda I$, with $\lambda$ the probe displacement from the target. $\lambda I$ is a scalar multiple of the identity and commutes with $H_T$, so the shifted propagator factorises,
\begin{equation*}
e^{-i t (H_T + \lambda I)} = e^{-i t \lambda} \, e^{-i t H_T} .
\end{equation*}
The probe enters the record only through the scalar phase $e^{-i t \lambda}$. A single banked record therefore serves every probe value in classical post-processing - this factorisation is the derivation behind the classical probe dial.

% NA1-p25
The first resource inversion converts window geometry and weight ratio into the filter strength the protocol must reach. The two window bounds combine into a ratio bound, the accuracy demand caps the ratio at $\eta$, and solving the capped bound yields the filter strength,
\begin{align*}
\frac{Z_\tau^{\mathrm{out}}}{Z_\tau^{\mathrm{in}}}
&\le \frac{P_{\mathrm{out}}}{P_{\mathrm{in}}} \, e^{-\tau (\Delta^2 - r^2)} , \\
\frac{P_{\mathrm{out}}}{P_{\mathrm{in}}} \, e^{-\tau (\Delta^2 - r^2)} &\le \eta , \\
\tau &\ge \frac{\log(P_{\mathrm{out}}/P_{\mathrm{in}}) + \log(1/\eta)}{\Delta^2 - r^2} .
\end{align*}
$\tau$ stays polynomial exactly when both logarithms are $O(\log n)$ and the gap term $\Delta^2 - r^2$ is at least inverse polynomial in $n$. The polynomiality reading is conditional - it holds only when the weight ratio and the window gap satisfy the stated scalings, and those conditions are precisely what the body's obstruction analysis interrogates.

% NA1-p26
The tail bound is what makes the finite record sufficient. Truncating the quadrature at $\bar m$ retained pairs leaves the additive error
\begin{equation*}
\varepsilon_{\bar m} \le \sum_{k > \bar m} w_k f_k ,
\end{equation*}
and the Gaussian weights collapse super-exponentially once the grid covers the measured range, so the tail is controlled by the grid itself. With Nyquist spacing $\Delta t = \pi / (\sqrt{8} \, \|H\|)$ and the aperture reach $t_{\max} = \sqrt{\tau \ln(1/\varepsilon_w)}$~\eqref{res:aperture-resolution} in the appendix width convention, the retained budget is
\begin{widetext}
\begin{equation*}
\bar m = \frac{t_{\max}}{\Delta t}
= \frac{\sqrt{8}}{\pi} \, \|H\| \sqrt{\tau \ln(1/\varepsilon_w)}
\propto \sqrt{s \, \ln(1/\varepsilon_w)} , \qquad s = \tau \|H\|^2 ,
\end{equation*}
\end{widetext}
with $s$ the effective spectral bandwidth - the $\bar m$ entry of the resource triple, at the resolution the record section quotes. The collapse constant of the tail is representation dependent, and the budget is a scaling statement rather than a sharp constant.

% NA1-p27
The shot inversion completes the resource triple by pricing the statistical error of the same record. The estimator error $\varepsilon_K = c K^{-1/2}$~\eqref{def:shot-budget-K}, with $c$ the coefficient norm of the reconstruction estimator, must sit below $\eta$ times the in-window floor,
\begin{align*}
c \, K^{-1/2} &\le \eta \, P_{\mathrm{in}} \, e^{-\tau r^2} , \\
K &\ge c^2 \, \eta^{-2} \, P_{\mathrm{in}}^{-2} \, e^{2 \tau r^2} , \\
K &\ge c^2 \, \eta^{-2} \, P_{\mathrm{in}}^{-2} \qquad (r = 0) ,
\end{align*}
the last line the centred probe. The demand is stated against the in-window lower bound $P_{\mathrm{in}} \, e^{-\tau r^2}$ rather than the exact $Z_\tau^{\mathrm{in}}$, so the bound is conservative. $K$ stays polynomial exactly when $1/P_{\mathrm{in}}$ and the constant $c$ are both polynomial in $n$ - and whether $P_{\mathrm{in}}$ is large enough is exactly the question the cluster analysis of the next subsection addresses.

\subsection*{Cluster block reduction and the second-order target spread}

% NA1-p28
The resource inversions stay polynomial only if the perturbed spectrum keeps enough weight in the window, and the perturbative content that decides it is the cluster effective Hamiltonian of the body,
\begin{widetext}
\begin{equation*}
H_{\mathrm{eff}}^{(S)}(z) = S P_S + P_S V P_S + \Sigma_S(z) , \qquad
\Sigma_S(z) = P_S V Q_S \left( z - Q_S H Q_S \right)^{-1} Q_S V P_S ,
\end{equation*}
\end{widetext}
defined for $z$ in the cluster window $|z - S| < \Gamma_S/2$, derived here by Riesz projection and similarity, then followed through the perturbed eigenvectors to the overlap weights and the second-order spread of the target cluster. The body's effective-Hamiltonian display and the multiplicity-fingerprint forward model draw on its outputs.

% NA1-p29
Let $P_S = \sum_{x : s(x) = S} |x\rangle\langle x|$ project onto the subset-sum-$S$ eigenspace of $H_0$, with complement $Q_S = I - P_S$ and isolation gap $\Gamma_S = \mathrm{dist}(S, \mathrm{spec}(Q_S H_0 Q_S))$, at least $1$ for integer weights, and let the perturbation $V$ be bounded, Hermitian and non-commuting with at least one $Z_i$, in the local two-body form $V = \sum_{a < b} V_{ab}$. For $H = H_0 + V$ with $\|V\| < \Gamma_S/4$ the Riesz projector onto the perturbed cluster is
\begin{equation*}
\Pi_S = \frac{1}{2\pi i} \oint_{|z - S| = \Gamma_S/2} (z - H)^{-1} \, dz .
\end{equation*}
Weyl stability moves every perturbed eigenvalue by at most $\|V\| < \Gamma_S/4$, so the cluster eigenvalues stay inside the contour \cite{bhatia1997matrix}, and the isolation gap keeps the remaining spectrum at distance at least $3\Gamma_S/4$ from $S$, outside it. The projector is therefore well defined \cite{kato1995perturbation}. The construction needs only the norm condition and Hermiticity - no locality assumption enters until the spread estimate that closes the subsection.

% NA1-p30
The projector pair reduces the Hamiltonian to the cluster block by a similarity. The direct rotation built from the two projectors intertwines them, and conjugating $H$ and restricting to the cluster yields the effective Hamiltonian,
\begin{align*}
U_S &= \Pi_S P_S + (I - \Pi_S)(I - P_S) , \\
\Pi_S &= U_S P_S U_S^{-1} , \\
H_{\mathrm{eff}}^{(S)} &= P_S U_S^{-1} H U_S P_S = S P_S + P_S V P_S + \Sigma_S ,
\end{align*}
with the three terms arising as the zeroth, first and second order of the block decomposition \cite{kato1995perturbation}. The final term is the Schur complement of the outer block - $\Sigma_S$ collects the second-order virtual transitions through $Q_S H Q_S$ \cite{horn2013matrix}. It is evaluated at spectral parameters in the isolation annulus, where the resolvent of $Q_S H Q_S$ exists, and its norm is second order, $\|\Sigma_S\| = O(\|V\|^2 / \Gamma_S)$.

% NA1-p31
The block reduction becomes amplitude statements once the perturbed eigenvectors are decomposed over the cluster strings. Writing $C(S) = \{ x : s(x) = S \}$ for the $g(S)$ cluster strings, a perturbed cluster eigenvector reads
\begin{equation*}
|E_{S,\alpha}^{\mathrm{pert}}\rangle = \sum_{x \in C(S)} u^{\alpha}_x \, |x\rangle , \qquad
\sum_{x \in C(S)} |u^{\alpha}_x|^2 = 1 ,
\end{equation*}
with $\{u^\alpha\}_\alpha$ the eigenvectors of the internal block $P_S V P_S$ in the cluster basis and the support staying inside the cluster up to corrections of order $\|V\|/\Gamma_S$. The internal block also fixes the energies - the first-order intra-cluster splitting equals the spectrum of $P_S V P_S$, the degeneracy of $H_0$ lifted by exactly its eigenvalues $\nu_\alpha$. Generic $V$ gives all distinct $\nu_\alpha$, while structured $V$ may preserve sub-degeneracies. The overlap computation uses this decomposition string by string.

% NA1-p32
Pairing the amplitudes against the uniform expansion of the plus state gives the overlap in closed form,
\begin{align*}
\langle E_{S,\alpha}^{\mathrm{pert}} | \psi_+ \rangle
&= \sum_{x \in C(S)} u^{\alpha}_x \, \langle x | \psi_+ \rangle
= 2^{-n/2} \sum_{x \in C(S)} u^{\alpha}_x , \\
|\langle E_{S,\alpha}^{\mathrm{pert}} | \psi_+ \rangle|^2
&= 2^{-n} \, \Big| \sum_{x \in C(S)} u^{\alpha}_x \Big|^2 .
\end{align*}
Every cluster string contributes the same coefficient $2^{-n/2}$, so the amplitudes interfere before they are squared. The weight is a coherent sum of cluster-restricted eigenvector amplitudes rather than an incoherent count of cluster states - this coherent structure is the mechanism behind the difference between perturbed weights and classical multiplicities, which is why the sum is displayed rather than asserted.

% NA1-p33
The coherent sums give content to the perturbed spectral weights,
\begin{equation*}
p_j^{\mathrm{perturbed}} = |\langle E_j^{\mathrm{pert}} | \psi \rangle|^2 ,
\end{equation*}
the spectral weights of the initial state on the perturbed eigenbasis. For the broad-support state they are the coherent sums just computed and no longer the classical multiplicities $g(T)/2^n$. The two coincide only when the internal-block eigenvectors align with single strings - coherence generally moves weight between cluster eigenvectors. The body cites this definition wherever perturbed weights replace counts.

% NA1-p34
The target cluster closes the appendix. For generic distinct integer weights and generic $k$-local $V$ the internal block on the target subspace vanishes,
\begin{equation*}
P_S V P_S = 0 \qquad (S = T ,\ \text{generic case}) ,
\end{equation*}
so first-order splitting is absent and the leading spread is second order, through the Schur complement \cite{horn2013matrix, bhatia1997matrix}. The target cluster spreads as
\begin{equation*}
w_{\mathrm{target}} = O\!\left( \frac{\|V\|^2 \, n^{k/2}}{\Gamma_S} \right) ,
\end{equation*}
with the $g(T)$ sub-eigenvalues distributed on a Wigner semicircle at $\{ T + \nu_\alpha \}$, the displacements now supplied by the second-order block. The first-order vanishing is a genericity statement about weights and perturbation rather than an identity, and the semicircle is the distributional summary of the sub-eigenvalues rather than a per-instance guarantee. This spread supplies the window support whose bandwidth the fingerprint forward model and the polynomial shot-budget argument draw on.

% PLAN:NA2
% NODES: res:window-verdict res:retention-frontier-nyquist irf:window-entangling-but-localized res:j-dial-fingerprint-scaling res:degeneracy-cartography-single-family res:integer-decoder def:erasure-rule def:integer-snap-and-trial-median irf:n14-window-chi-flat-through-n20 res:planted-trade-family thm:staircase-cumulant-identities thm:response-identities-s irf:tau-is-a-classical-dial
\section{Data availability and machine manifest}\label{app:NA2}

% NA2-p1
Every numerical result derives from classical simulation - the manuscript contains no quantum hardware data anywhere. The complete generating code, configuration file and banked data ship with the manuscript source. Three versioned banks hold every file a caption cites - the v6 exact-diagonalisation and exact-spectral-sum bank, the v7 matrix-product bank and the v8 continuation bank - each with machine-readable run manifests. Caption data-source identifiers ED-n, SYN-n, TEBD-n and XS-n are the only data pointers in prose and captions, and each resolves through the manifest table of this appendix. Repository paths and commit hashes appear only inside that table.

% NA2-p2
The manifest is Table~\ref{tab:manifest} - a row set per caption identifier carrying the six manifest fields: generator script, engine, configuration, seed lineage, banked data file, and commit. The prefix records how each file was computed: ED-n is exact diagonalisation at $n \le 12$, SYN-n synthetic finite-shot sampling at stated budget $K$, TEBD-n matrix-product evolution at $n \le 20$ and XS-n exact spectral sums evaluated numerically. The commit column is the abbreviated hash of the last commit touching the file - a per-file repository hash, not a release tag. Running the named script under the named configuration and seed lineage regenerates the named file at that commit. The v8 manifests also record interpreter and library versions - regeneration is reproducible up to platform floating point.

% NA2-p3
Four identifiers resolve to the v6 bank at commit \texttt{ec5b557}, under configuration \texttt{config\_v3.yaml} with master seed 20260701 and instance lineage $20260701 + 1000n + \mathrm{instance}$. ED-N1r maps to \texttt{N1r\_retention\_frontier.csv} - 75 rows from the \texttt{run\_N1r} experiment of the engine module \texttt{numerics\_v3.py}, the single-instance retention frontier behind panel (c) of the envelope figure. ED-N11 maps to \texttt{N11\_entanglement\_growth.csv} - 432 rows of half-cut entanglement and required bond dimension for the $V_{\mathrm{XXZ}}$ family. ED-N13 maps to \texttt{N13proto\_fingerprint.csv} from the sweep script \texttt{n13proto\_fingerprint\_jsweep.py} over the same engine module. XS-N12 maps to \texttt{N12\_identities.csv} from the N12 cartography experiment, whose companions \texttt{N12\_steps.csv}, \texttt{N12\_gchain.csv}, \texttt{N12\_ledger.csv} and \texttt{N12\_tilts.csv} bank the staircase steps, chain records and tilt responses that later rows reuse. The bank logs each experiment in \texttt{run\_manifest.jsonl} with per-experiment target-check files. The ED-N1r file holds one instance per system size - its instance-band extension is the deferred item recorded at the close of this manifest.

% NA2-p4
Work item WK-051 extends the fingerprint identifier with two v8 files at commit \texttt{44291de} - \texttt{N13b\_fingerprint\_distribution.csv} and \texttt{N13b\_separation.csv}, with \texttt{N13b\_manifest.json} banked alongside. The rows name the engine script \texttt{n13proto\_fingerprint\_jsweep.py} over the engine module \texttt{numerics\_v3.py} under \texttt{config\_v3.yaml}, with instance seed lineage $20260701 + 1000n + \mathrm{instance}$ and noise master seed 20260709. The bank holds per-target response distributions for $n \in \{8, 10\}$ with four instances per $n$ in the $V_{\mathrm{XXZ}}$ family at $\tau = 4$, both null channels, and shot budgets $K \in \{10^{3}, 10^{4}, 10^{5}\}$ with 250 realisations per cell and bootstrap confidence intervals from 300 resamples. These finite-shot channels are synthetic sampling from exact overlaps at the stated budgets, not hardware counts. The files supply the separation and distribution panels of the fingerprint-response figure, and the caption cites them by path under the parent fingerprint identifier rather than a separate prefix.

% NA2-p5
SYN-N12c maps to \texttt{v8/N12c\_decoder\_statistics.csv} at commit \texttt{22d7e2c}, produced by \texttt{n12c\_decoder\_statistics.py} as pure post-processing of the banked \texttt{v6/N12\_steps.csv} record of 33000 rows - no quantum engine is re-run. The run draws shot budgets $K \in \{0, 10^{3}, 10^{4}, 10^{5}\}$ with eight shot realisations per finite budget and noise seed 20260709, reports the single-trial and trial-median methods under the erasure variant and a separate ratio-one imputation variant, covers 24 instances and 264 targets, and attaches 500 bootstrap resamples. The erasure rule applied in post-processing is the same function the decoder itself applies - erased steps stay erased and are excluded from conditional accuracy, and the imputation variant is a prior-dependent curve reported separately, not merged into the erasure results. The file supplies the decoder-performance figure.

% NA2-p6
Three row sets serve the window-scaling identifier TEBD-N14. The v7 file \texttt{N14\_tebd\_scaling.csv} at commit \texttt{ec5b557} comes from \texttt{n14\_tebd\_hardness\_scaling.py}, reading \texttt{config\_v3.yaml}, instances seeded as $20260701 + 1000n + \mathrm{instance}$. Work item WK-053 banks \texttt{v8/N14c\_ed\_tebd\_overlap.csv} at commit \texttt{44291de} with \texttt{N14c\_manifest.json}, running both engines unchanged - exact diagonalisation through \texttt{numerics\_v3.SubsetSumInstance}, matrix-product evolution \cite{schollwock2011density} through \texttt{n14\_tebd\_hardness\_scaling.evolve} with second-order Trotter bond gates - on shared $n \in \{8, 10, 12\}$, two instances, common couplings $J \in \{0.05, 0.1, 0.2, 0.5\}$. Both run to the common maximum time $2\sqrt{\tau_{\max}\log(1/\epsilon_{w})} \approx 38.4$ ($\tau_{\max} = 16$, $\epsilon_{w} = 10^{-10}$), per-gate truncation $10^{-12}$, bond dimension at the $10^{-4}$ discarded-weight convention - identical for both methods - and the cap of 128 never reached. The bank holds 288 agreement rows with largest entropy discrepancy $2.7 \times 10^{-4}$ - direct agreement shown on $n \le 12$ only, sizes above the shared range supported by that shown agreement, not by comparison at those sizes. Work item WK-054 banks \texttt{v8/N14d\_instances.csv} at commit \texttt{44291de} with \texttt{N14d\_manifest.json} from \texttt{n14d\_instances.py}, driving the same \texttt{evolve} routine and seed lineage, four instances per $n \in \{8, 10, 12, 14, 16\}$ at primary coupling $J = 0.2$ on a twelve-point time grid from the banked N11 record, supplying the cost panel's instance spread.

% NA2-p7
ED/TEBD-N15b maps to three v8 files at commit \texttt{22d7e2c} - \texttt{N15b\_planted\_panels\_spread.csv}, \texttt{N15b\_planted\_panels\_chi.csv} and \texttt{N15b\_planted\_panels\_entropy.csv} - banked by work item WK-055 with \texttt{N15b\_manifest.json} from the panel script \texttt{n15b\_planted\_panels.py}. The manifest fixes linear chain ordering with block sites first, Vidal canonical form with second-order Trotter bond gates, per-gate truncation $10^{-12}$, bond-dimension cap 128, time step 0.005, the $10^{-4}$ discarded-weight convention and the coupling convention $J = 0.2/L$. Provenance is split by panel: the in-window spread panel is recomputed by exact diagonalisation at $n \in \{8, 12\}$, with $n \ge 16$ marked infeasible for exact diagonalisation, while the bond-dimension and entropy panels reuse \texttt{v7/N15\_planted\_growth.csv} verbatim from \texttt{n15\_planted\_trade\_demo.py} at commit \texttt{ec5b557}, nothing recomputed. Bond-dimension growth for the planted family is consistent with exponential over the tested sizes \cite{eisert2010area} and is recorded as a guide rather than an established law, matching the figure caption.

% NA2-p8
Figure~\ref{fig:N12-identities} evaluates one banked exact-overlap record - the XS-N12 file of the v6 rows - through two reconstruction routes under shared code. It plots the median relative error, per dial $\tau$ and perturbation family $V = 0$, $V_{\mathrm{XX}}$ and $V_{\mathrm{XXZ}}$, of the staircase derivative identity $\mathrm{d}H/\mathrm{d}\lambda = -2\tau\sigma_{\tau}^{2}$ and of the agreement between the derivative route and the coefficient-insertion route. The derivative is a second-order central difference on a probe-shift grid of spacing $5 \times 10^{-4}$, and the quadrature retains the per-row node count recorded in the data file. The identity error sits at the grid-differencing floor near $10^{-5}$, the route agreement at the quadrature floor near $10^{-12}$, and the $\tau$-dial subset-exactness error is exactly zero on every row - stated here rather than plotted on a logarithmic axis. The figure supports one statement: the implementation is internally consistent and was checked against direct spectral sums. It is not evidence for the algebraic identities, which are proved, and involves no finite-shot noise.

% NA2-p9
One numeric remains deferred. The retention frontier of ED-N1r is banked at v6 as a single instance per system size, and its instance-band extension awaits its run - instance-to-instance spread of the frontier is not yet established, and the envelope figure carries the matching single-instance qualification in its caption. The deferral concerns the band extension alone - the banked frontier file is complete and resolves normally. Every other caption identifier resolves through the manifest to a generator script, an engine, a configuration, a seed lineage, a banked data file and a commit.

% table environment generated for AO3 (canonised-run): the AppB machine manifest as a
% real table, one row per caption identifier, replacing the prose catalogue. Seven
% columns (caption identifier + the six manifest fields). Rebuilt at full size: each
% cell is a fixed-width parbox and script/file paths are made breakable (underscores
% and slashes carry a break opportunity) so long identifiers wrap inside their column
% width rather than being shrunk to fit; reduced \tabcolsep, banked-file lists comma-separated.
\begin{table*}[tb]
\setlength{\tabcolsep}{3.5pt}
% breakable path tokens, scoped to this table: underscore and slash keep their glyph
% but gain a zero-width break opportunity so paths wrap at the fixed column width.
\renewcommand{\_}{\textunderscore\allowbreak}
\caption{Machine manifest. Each caption data-source identifier resolves to one row: the generator script, the engine module, the configuration, the seed lineage, the banked data file(s) and the abbreviated per-file commit hash. Running the named script under the named configuration and seed lineage regenerates the named file at that commit. Prefixes record the computation: ED exact diagonalisation ($n\le12$), SYN synthetic finite-shot sampling at budget $K$, TEBD matrix-product evolution ($n\le20$), XS exact spectral sums. Bank version and interpreter/library versions are recorded per bank - regeneration is reproducible up to platform floating point.}
\label{tab:manifest}
\begin{ruledtabular}
\begin{tabular}{lllllll}
\parbox[t]{0.08\textwidth}{\raggedright\textbf{Identifier}\strut} & \parbox[t]{0.15\textwidth}{\raggedright\textbf{Generator script}\strut} & \parbox[t]{0.10\textwidth}{\raggedright\textbf{Engine}\strut} & \parbox[t]{0.09\textwidth}{\raggedright\textbf{Config}\strut} & \parbox[t]{0.15\textwidth}{\raggedright\textbf{Seed lineage}\strut} & \parbox[t]{0.20\textwidth}{\raggedright\textbf{Banked file(s)}\strut} & \textbf{Commit} \\ \hline
\parbox[t]{0.08\textwidth}{\raggedright ED-N1r\strut} & \parbox[t]{0.15\textwidth}{\raggedright run\_N1r\strut} & \parbox[t]{0.10\textwidth}{\raggedright numerics\_v3.py\strut} & \parbox[t]{0.09\textwidth}{\raggedright config\_v3.yaml\strut} & \parbox[t]{0.15\textwidth}{\raggedright 20260701 + 1000n + instance\strut} & \parbox[t]{0.20\textwidth}{\raggedright N1r\_retention\_frontier.csv\strut} & ec5b557 \\[4pt]
\parbox[t]{0.08\textwidth}{\raggedright ED-N11\strut} & \parbox[t]{0.15\textwidth}{\raggedright run\_N11\strut} & \parbox[t]{0.10\textwidth}{\raggedright numerics\_v3.py\strut} & \parbox[t]{0.09\textwidth}{\raggedright config\_v3.yaml\strut} & \parbox[t]{0.15\textwidth}{\raggedright 20260701 + 1000n + instance\strut} & \parbox[t]{0.20\textwidth}{\raggedright N11\_entanglement\_growth.csv\strut} & ec5b557 \\[4pt]
\parbox[t]{0.08\textwidth}{\raggedright ED-N13\strut} & \parbox[t]{0.15\textwidth}{\raggedright n13proto\_fingerprint\_jsweep.py\strut} & \parbox[t]{0.10\textwidth}{\raggedright numerics\_v3.py\strut} & \parbox[t]{0.09\textwidth}{\raggedright config\_v3.yaml\strut} & \parbox[t]{0.15\textwidth}{\raggedright 20260701 + 1000n + instance\strut} & \parbox[t]{0.20\textwidth}{\raggedright N13proto\_fingerprint.csv\strut} & ec5b557 \\[4pt]
\parbox[t]{0.08\textwidth}{\raggedright XS-N12\strut} & \parbox[t]{0.15\textwidth}{\raggedright run\_N12\strut} & \parbox[t]{0.10\textwidth}{\raggedright numerics\_v3.py\strut} & \parbox[t]{0.09\textwidth}{\raggedright config\_v3.yaml\strut} & \parbox[t]{0.15\textwidth}{\raggedright 20260701 + 1000n + instance\strut} & \parbox[t]{0.20\textwidth}{\raggedright N12\_identities.csv, N12\_steps.csv, N12\_gchain.csv, N12\_ledger.csv, N12\_tilts.csv\strut} & ec5b557 \\[4pt]
\parbox[t]{0.08\textwidth}{\raggedright N13b (WK-051)\strut} & \parbox[t]{0.15\textwidth}{\raggedright n13proto\_fingerprint\_jsweep.py\strut} & \parbox[t]{0.10\textwidth}{\raggedright numerics\_v3.py\strut} & \parbox[t]{0.09\textwidth}{\raggedright config\_v3.yaml\strut} & \parbox[t]{0.15\textwidth}{\raggedright 20260701 + 1000n + instance, noise 20260709\strut} & \parbox[t]{0.20\textwidth}{\raggedright N13b\_fingerprint\_distribution.csv, N13b\_separation.csv\strut} & 44291de \\[4pt]
\parbox[t]{0.08\textwidth}{\raggedright SYN-N12c\strut} & \parbox[t]{0.15\textwidth}{\raggedright n12c\_decoder\_statistics.py\strut} & \parbox[t]{0.10\textwidth}{\raggedright post-proc.\ of v6/N12\_steps.csv\strut} & \parbox[t]{0.09\textwidth}{\raggedright config\_v3.yaml\strut} & \parbox[t]{0.15\textwidth}{\raggedright noise 20260709\strut} & \parbox[t]{0.20\textwidth}{\raggedright v8/N12c\_decoder\_statistics.csv\strut} & 22d7e2c \\[4pt]
\parbox[t]{0.08\textwidth}{\raggedright TEBD-N14\strut} & \parbox[t]{0.15\textwidth}{\raggedright n14\_tebd\_hardness\_scaling.py\strut} & \parbox[t]{0.10\textwidth}{\raggedright TEBD (evolve)\strut} & \parbox[t]{0.09\textwidth}{\raggedright config\_v3.yaml\strut} & \parbox[t]{0.15\textwidth}{\raggedright 20260701 + 1000n + instance\strut} & \parbox[t]{0.20\textwidth}{\raggedright N14\_tebd\_scaling.csv\strut} & ec5b557 \\[4pt]
\parbox[t]{0.08\textwidth}{\raggedright N14c (WK-053)\strut} & \parbox[t]{0.15\textwidth}{\raggedright n14\_tebd\_hardness\_scaling.py\strut} & \parbox[t]{0.10\textwidth}{\raggedright numerics\_v3.py + TEBD\strut} & \parbox[t]{0.09\textwidth}{\raggedright config\_v3.yaml\strut} & \parbox[t]{0.15\textwidth}{\raggedright 20260701 + 1000n + instance\strut} & \parbox[t]{0.20\textwidth}{\raggedright v8/N14c\_ed\_tebd\_overlap.csv\strut} & 44291de \\[4pt]
\parbox[t]{0.08\textwidth}{\raggedright N14d (WK-054)\strut} & \parbox[t]{0.15\textwidth}{\raggedright n14d\_instances.py\strut} & \parbox[t]{0.10\textwidth}{\raggedright TEBD (evolve)\strut} & \parbox[t]{0.09\textwidth}{\raggedright config\_v3.yaml\strut} & \parbox[t]{0.15\textwidth}{\raggedright 20260701 + 1000n + instance\strut} & \parbox[t]{0.20\textwidth}{\raggedright v8/N14d\_instances.csv\strut} & 44291de \\[4pt]
\parbox[t]{0.08\textwidth}{\raggedright ED/TEBD-N15b\strut} & \parbox[t]{0.15\textwidth}{\raggedright n15b\_planted\_panels.py\strut} & \parbox[t]{0.10\textwidth}{\raggedright ED + TEBD\strut} & \parbox[t]{0.09\textwidth}{\raggedright config\_v3.yaml\strut} & \parbox[t]{0.15\textwidth}{\raggedright 20260701 + 1000n + instance\strut} & \parbox[t]{0.20\textwidth}{\raggedright N15b\_planted\_panels\_spread.csv, \_chi.csv, \_entropy.csv\strut} & 22d7e2c \\
\end{tabular}
\end{ruledtabular}
\end{table*}

% figure environment generated by tex_backend from the figure spec
% spec: fig:N12-identities (status: banked)
\begin{figure}[tb]
    \centering
    \includegraphics[width=\linewidth]{figures/N12_identities.pdf}
    \caption{Median relative error of the staircase derivative identity $dH/d\lambda = -2\tau\sigma^{2}$ and of the agreement between the derivative reconstruction route and the coefficient-insertion route, plotted per dial $\tau$ for the families $V = 0$, $V_{\mathrm{XX}}$ and $V_{\mathrm{XXZ}}$ on one banked exact-overlap record (data source XS-N12). The derivative is a second-order central difference on a probe-shift grid of spacing $5 \times 10^{-4}$, the reconstruction quadrature retains the per-row node count recorded in the data file, and both routes evaluate the same banked record through the same reconstruction code path. The identity error sits at the grid-differencing floor near $10^{-5}$ and the route agreement at the quadrature floor near $10^{-12}$, while the $\tau$-dial subset-exactness error is exactly zero on every row and is therefore stated in text rather than plotted on a logarithmic axis. The plotted errors are an internal-consistency check against direct spectral sums rather than evidence for the algebraic identities themselves, which are proved, and involve no finite-shot noise.}
    \label{fig:N12-identities}
\end{figure}


% The disclosure-crosswalk appendix (NA3.tex + tables_tex/tab_crosswalk.tex) was
% RETIRED at r6 by git rm. The companion invention-disclosure document now carries
% the crosswalk as its own table and is its home - git history preserves the
% retired section. A one-sentence pointer closes the back matter (the
% data-availability appendix above).
\noindent The disclosure crosswalk mapping each supported statement to its technical backing is carried by the companion invention-disclosure document.

\bibliography{references}

\end{document}
